\selectlanguage{ngerman}%
\frenchspacing%
\begin{ledgroupsized}[r]{120mm}%
\footnotesize%
\pstart%
\noindent\textbf{Überlieferung:}%
\pend%
\end{ledgroupsized}%
\begin{ledgroupsized}[r]{114mm}%
\footnotesize%
\pstart%
\parindent -6mm%
\makebox[6mm][l]{\textit{L}}%
Konzept:
LH~XXXVII~3 Bl.\:123--124.
Ein Bogen 2\textsuperscript{o};
Wasserzeichen auf Bl.\:\href{\Scan{LH_37_03_124wz}}{124};
Ränder teilweise ausgefranst mit geringfügigem Textverlust
(Papiererhaltungsmaßnahmen);
nachträgliche Textänderungen im oberen Drittel von Bl.\:\href{\Scan{LH_37_03_123r}}{123\,r\textsuperscript{o}}\! und in der Mitte von Bl.\:\href{\Scan{LH_37_03_124v}}{124\,v\textsuperscript{o}}\! mit verblasster Tinte verfasst.
Vier ganzseitig beschriebene Seiten
mit zahlreichen Verbesserungen und Nebenrechnungen;
Textfolge grundsätzlich nach Blattzählung,
wobei der Textabschnitt \refpassage{LH_37_03_123-124_123r+124r_zweitestufe-1}{LH_37_03_123-124_123r+124r_zweitestufe-2}
teils auf Bl.\:\href{\Scan{LH_37_03_123r}}{123\,r\textsuperscript{o}}\! und
teils auf Bl.\:\href{\Scan{LH_37_03_124r}}{124\,r\textsuperscript{o}}\! verfasst ist
(ein Kustos verbindet beide Textteile).
Auf Bl.\:\href{\Scan{LH_37_03_124r}}{124\,r\textsuperscript{o}}\! im oberen Seitenbereich
ein zum Teil gestr. Vermerk von Leibnizens Hand:
\textit{%
\sout{%
Ich solte dafür halten, daß Schediasma de angulo incidentiae et reflexionis aequali sey auß zu lassen weilen, in der that darinn keine ratio dieser%
}
Putem omittendam esse dissertationem de angulo incidentiae et reflexionis aequali, quia revera in ea ostenditur quidem corpus oblique incidens non debere recedere e linea qua venit, sed non ostenditur quod angulus lineae recessus sit aequalis angulo lineae adventus. Quod aliunde satis explicatum habetur.%
}
Die in N.\:\ref{37109+60209}, \refpassage{MB_1710_123-128+LH_37_03_120-122_121r_berechng}{MB_1710_123-128+LH_37_03_120-122_121r_berechng}\!\,ff. angestellte Berechnung knüpft an das lediglich in N.\:\ref{60210}, Bl.\:\href{\Scan{LH_37_03_123r}}{123\,r\textsuperscript{o}}\! (S.\:\pageref{LH_37_03_123-124_123r_Fig1})
überlieferte Diagramm \lbrack\textit{Fig.\:1}\rbrack\ an.
\pend%
\end{ledgroupsized}%
\selectlanguage{latin}%
\frenchspacing%
\vspace{8mm}%
\normalsize%
\pstart%
\noindent%
\lbrack%
\href{\Scan{LH_37_03_123r}}%
{123\,r\textsuperscript{o}}%
\!.
\textit{Nachfolgend kleingedruckter Text in L gestrichen und durch den überschriebenen Textabschnitt \refpassage{LH_37_03_123-124_123r_ersetzung-1}{LH_37_03_123-124_123r_ersetzung-2} ersetzt:}%
\rbrack%
\pend%
\vspace{1.0em}%
\footnotesize%
\pstart%
\noindent%
\centering%
Annotatio%
\indexS{annotatio}
\edlabel{LH_37_03_123-124_123r_anfang-1}%
G.G.L.%
\indexN{\textls{Leibniz} (Leibnitius, GGL), Gottfried Wilhelm 1646\textendash1716}%
\edtext{%
}{%
\xxref{LH_37_03_123-124_123r_anfang-1}{LH_37_03_123-124_123r_anfang-2}%
\lemma{G.G.L.}%
\Bfootnote{\lbrack/\rbrack\ %
\textit{(1)}~Difficultatis aliquid in hac explicatione reperio. %
\textit{(2)}~In hac %
\lbrack...\rbrack\ immutandum putem.%
~\textit{L}%
}}%
\pend%
\vspace{0.5em}%
\pstart%
\noindent%
In
\edtext{%
hac Explicatione%
\indexS{explicatio}%
}{%
\lemma{hac Explicatione}%
\Cfootnote{%
Vgl. E.\;\textsc{Chauvin}, \glqq Nova circa vapores hypothesis\grqq, \textit{MB} 1710, S.\:120--122.\cite{01415}
Siehe über die Ümstände der Entstehung von N.\:\ref{60210}\! die einführende Erläuterung, S.\:\pageref{elevatio_intro}\,f.%
}}
aliquid supplendum, et aliquid immutandum putem.%
\edlabel{LH_37_03_123-124_123r_anfang-2}
Vult doctissimus
\edtext{%
autor%
\indexS{autor doctus}
in ascensu vaporum%
\indexS{ascensus vaporis}
bullas%
\indexS{bulla vaporis}%
}{%
\lemma{autor}%
\Bfootnote{%
\textit{(1)}~bull\textlangle\textendash\textrangle\ %
\textit{(2)}~bullas %
\textit{(3)}~in ascensu vaporum bullas%
~\textit{L}%
}}
ex aquae pellicula%
\indexS{pellicula aquea bullarume}
et aere incluso rarefacto%
\indexS{aer inclusus bullae}%
\indexS{aer rarefactus}
constantes
pelli sursum in aerem communi,%
\indexS{bulla pulsa sursum}%
\indexS{aer communis}
non per pondus aeris ambientis,%
\indexS{pondus aeris ambientis}%
\indexS{aer ambiens}
sed per vim ejus
\edtext{%
Elasticam.%
\indexS{vis elastica aeris ambientis}
Sed \textls{primo} vis Elastica aeris externi%
\indexS{vis elastica aeris externi}%
\indexS{aer externus}
aut eadem erit quae inclusi,%
\indexS{vis elastica aeris inclusi}
aut si major sit
comprimere conabitur aerem bullae inclusum,%
\indexS{aer inclusus bullae}
elastro debiliore praeditum;%
\indexS{elastrum aeris}
atque illa bullam vel reddet minorem%
\indexS{bulla reddita minor}
vel dirumpet.%
\indexS{bulla dirupta}
Desideratur igitur%
}{%
\lemma{Elasticam}%
\Bfootnote{%
\textit{(1)}~. Hic duo desiderantur, primum causa cur aer %
\textit{(2)}~. Sed vis Elastica aeris %
\textit{(3)}~, atque ita bullam vel reddet %
\textit{(a)}~arc %
\textit{(b)}~minorem vel dirumpet %
\textit{(4)}~Sed \textls{primo} vis Elastica aeris %
\textit{(a)}~comprimere conabitur %
\textit{(b)}~externi %
\textit{(aa)}~non %
\textit{(bb)}~potius %
\textit{(cc)}~spatio %
\textit{(dd)}~\textlangle sur\textrangle\ %
\textit{(ee)}~eadem %
\textit{(ff)}~aut eadem %
\lbrack...\rbrack\ aut si %
\textit{(aaa)}~diversa %
\textit{(bbb)}~major sit comprimere %
\textit{(aaaa)}~nitetur %
\textit{(bbbb)}~conabitur aerem %
\lbrack...\rbrack\ debiliore praeditum; %
\textit{(aaaaa)}~explicanda ergo ratio erat %
\textit{(bbbbb)}~atque illa %
\lbrack...\rbrack\ Desideratur igitur%
~\textit{L}%
}}
\lbrack%
\textit{der Text wird auf \refpassage{LH_37_03_123-124_123r_ersetzung-2}{LH_37_03_123-124_123r_ersetzung-2} fortgesetzt.}%
\rbrack%
\newpage%
\pend%
\vspace{1.5em}%
\normalsize%
\pstart%
\noindent%
\centering%
\edlabel{LH_37_03_123-124_123r_ersetzung-1}%
\edtext{%
Annotatio%
\indexS{annotatio}
ad Hypothesin%
\indexS{hypothesis de ascensu vaporis}
de Ascensu vaporum%
\indexS{vapor ascendens calore}%
\indexS{ascensus vaporis}%
}{%
\lemma{Annotatio}%
\Bfootnote{%
\textit{(1)}~de Vapore %
\textit{(2)}~ad Hypothesin de Ascensu vaporum%
~\textit{L}%
}}
\pend%
\vspace*{0.5em}%
\pstart%
\noindent%
Partem Vaporum ex%
\indexS{vapor ascendens calore}
\edtext{%
bullis esse,%
\indexS{bulla vaporis}
quae ex aquae pellicula%
\indexS{pellicula aquea bullarume}
et aere incluso rarefacto constent,%
\indexS{aer inclusus bullae}%
\indexS{aer rarefactus}
verisimilitudine non caret.%
\indexS{verisimilitudo}
Et tales bullae%
\indexS{bulla vaporis}
si habeant%
}{%
\lemma{bullis}%
\Bfootnote{%
\textit{(1)}~constare %
\textit{(2)}~esse, quae %
\lbrack...\rbrack\ rarefacto constent, %
\textit{(a)}~non %
\textit{(b)}~verisimilitudine non caret. Et %
\textit{(aa)}~cum %
\textit{(bb)}~tales bullae %
\textit{(aaa)}~cum \textlangle\textendash\textrangle t %
\textit{(bbb)}~si habeant%
~\textit{L}%
}}
minorem gravitatem specificam%
\indexS{gravitas specifica bullae}
quam aer ambiens,%
\indexS{aer ambiens}
ascendent,%
\indexS{bulla ascendens in aere}
alias
\edtext{%
secus,
nam vis Elastica aeris ambientis%
\indexS{vis elastica aeris ambientis}
cum aeque in omnes partes agat,
per se ad ascensum nil facit,%
\indexS{ascensus vaporis}
nec agere bullas minus elasticas%
\indexS{bulla elastica}
\lbrack aut\rbrack\ comprimere tentabit.%
}{%
\lemma{secus,}%
\Bfootnote{%
\textit{(1)}~vi %
\textit{(2)}~nam vis Elastica %
\textit{(a)}~cum sursum %
\textit{(b)}~aeris ambientis %
\lbrack...\rbrack\ partes agat, %
\textit{(aa)}~per se nihil ad ascensum facit, %
\textit{(bb)}~cum vi dela %
\textit{(cc)}~per se %
\lbrack...\rbrack\ facit, nec %
\textit{(aaa)}~tam sursum %
\textit{(bbb)}~agere bullas %
\textit{(aaaa)}~\textlangle, qua\textrangle\ %
\textit{(bbbb)}~minus elasticas \bars{quam \textit{ändert Hrsg.}} comprimere tentabit.%
~\textit{L}%
}}
Sed %
vis Elastica aeris externi%
\indexS{vis elastica aeris externi}
si
\edtext{%
m\lbrack aj\rbrack or%
}{%
\lemma{minor}%
\Bfootnote{%
\textit{L\;ändert Hrsg.}%
}}
sit
\edtext{%
quam inclusi in bulla,%
\indexS{vis elastica aeris inclusi}%
}{%
\lemma{quam}%
\Bfootnote{%
\textit{(1)}~\textlangle aer\textrangle\ %
\textit{(2)}~inclusi in bulla,%
~\textit{L}%
}}
id est si aer inclusus%
\indexS{aer inclusus bullae}
sit rarefactus,%
\indexS{aer rarefactus}
comprimere eum conabitur;%
\indexS{aer compressus}%
\indexS{compressio aeris}
atque ita bullam
vel reddet minorem
vel
\edlabel{LH_37_03_123-124_123r_dirumpetItaque-1}%
dirumpet.%
\indexS{bulla dirupta}%
\edtext{%
}{%
\xxref{LH_37_03_123-124_123r_dirumpetItaque-1}{LH_37_03_123-124_123r_dirumpetItaque-2}%
\lemma{dirumpet.}%
\Bfootnote{%
\textit{(1)}~Ita nihil aget, \textlangle aut\textrangle\ si major sit, comprimere %
\textit{(2)}~\lbrack/\rbrack\ Itaque desideratur \bars{igitur \textit{streicht Hrsg.}} %
\textit{(a)}~aliqua ratio, cur comprimi se %
\textit{(aa)}~ab %
\textit{(bb)}~\textlangle\textendash\textrangle\ %
\textit{(cc)}~aer inclusus non patiatur. Id poterit %
\bars{%
\textit{(1)}~fortasse conf %
\textit{(2)}~posterius in calorem conferri \textit{erg. u. gestr.}} %
oriri ab %
\textit{(aaa)}~aliquo %
\textit{(bbb)}~ejus calore, sed %
\textit{(aaaa)}~non apparet %
\textit{(bbbb)}~reddenda ratio %
\textit{(b)}~\bars{aliqua \textit{gestr.}} ratio cur aer %
\textit{(c)}~duplex ratio %
\textit{(d)}~duplicis effectus ratio, %
\textit{(aa)}~nec %
\textit{(bb)}~tum cur %
\lbrack...\rbrack\ etiam cur %
\textit{(aaa)}~\textlangle\textendash\textrangle\ 
\textit{(bbb)}~comprimi %
\textit{(aaaa)}~se ab %
\textit{(bbbb)}~se ab ambiente non patiatur. %
\textit{(aaaaa)}~Potest %
\textit{(bbbbb)}~Et potest sane haec resistentia %
\bars{peti \textit{erg.}} a calore %
\textit{(aaaaa-a)}~peti; sed quaeritur rursus %
\textit{(bbbbb-b)}~aeris inclusi; sed rursus quaeritur%
~\textit{L}%
}}
\pend%
\pstart%
Itaque desideratur%
\edlabel{LH_37_03_123-124_123r_ersetzung-2}
duplicis effectus ratio,%
\indexS{ratio effectus}%
\indexS{effectus duplex}
tum cur aer inclusus%
\indexS{aer inclusus bullae}
sit rarior ambiente,%
\indexS{aer rarior ambiente}
tum etiam cur comprimi se ab ambiente%
\indexS{aer compressus}
non patiatur.%
\indexS{aer resistens compressioni}
Et potest sane haec resistentia%
\indexS{resistentia bullae compressioni}
peti a calore aeris inclusi;%
\indexS{calor aeris inclusi}%
\indexS{aer inclusus bullae}
sed rursus quaeritur%
\edlabel{LH_37_03_123-124_123r_dirumpetItaque-2}
unde calor%
\indexS{calor aeris inclusi}
\edtext{%
sit%
}{%
\lemma{sit}%
\Bfootnote{%
\textit{erg.\;L}%
}}
in aere incluso %
potius
\edtext{%
quam in aere ambiente,%
\indexS{calor aeris ambientis}
atque adeo ratio simul%
\indexS{ratio effectus}
prioris effectus petitur%
\lbrack:\rbrack\
cur ille rarior isto?
An dicemus%
}{%
\lemma{quam}%
\Bfootnote{%
\textit{(1)}~in aere ambiente, nisi %
\textit{(2)}~in aere ambiente, atque adeo %
\textit{(a)}~cur %
\textit{(b)}~ratio simul prioris %
\textit{(aa)}~effectus quaeritur %
\textit{(bb)}~effectus petitur cur %
\textit{(aaa)}~aer inclusus %
\textit{(aaaa)}~sit ra %
\textit{(bbbb)}~plus %
\textit{(cccc)}~rarior sit ambiente %
\textit{(bbb)}~\bars{cur \textit{streicht Hrsg.}} ille rarior isto? An dicemus%
~\textit{L}%
}}
bullam a loco%
\indexS{locus calidus}%
\indexS{bulla vaporis}
\edtext{%
calido
(ubi forte aer%
}{%
\lemma{calido}%
\Bfootnote{%
\textit{(1)}~et %
\textit{(2)}~ubi aer %
\textit{(3)}~ubi forte aer %
\textit{(4)}~(ubi forte aer%
~\textit{L}%
}}
ambiens%
\indexS{aer ambiens}
et inclusus%
\indexS{aer inclusus bullae}
aeque rari erant)
ad locum%
\indexS{locus frigidus}
\edtext{%
frigidum %
nescio quo causu %
transvectam,%
\indexS{bulla transvecta a loco calido in frigidum}
calore suo conservato;%
\indexS{calor bullae}%
\indexS{calor conservatus}
\ledpb%
an ponemus initio aeque calidum fuisse
aerem externum%
\indexS{aer externus}
et aerem bullae inclusum%
\lbrack,\rbrack\
sed inclusum%
\indexS{aer inclusus bullae}
diutius suum calorem conservasse.%
\indexS{calor conservatus}
An potius dicemus
a calore aerem intra aquam produci%
\indexS{aer productus calore intra aquam}
uti aerem%
\indexS{aer natus ex aqua imposita igni}
ex aqua igni imposita%
\indexS{aqua imposita igni}
nasci constat;
aerem autem sic natum%
\indexS{aer natus ex aqua imposita igni}
et in bulla inclusum%
\indexS{aer inclusus bullae}
longe aere pristino ambiente%
\indexS{aer ambiens}
esse tenuiorem;%
\indexS{aer tenuior ambiente}
et quod vi elasticae%
\indexS{vis elastica aeris}
seu densitati deest,%
\indexS{densitas aeris}
vi dilatativa%
\indexS{vis dilatativa aeris}
vel explosiva%
\indexS{vis explosiva aeris}
suppleri,%
\indexS{vis suppleta}
ne scilicet ab ambiente comprimatur.%
\indexS{aer compressus}%
\indexS{aer ambiens}
Porro
quod calor efficit,%
\indexS{effectus caloris}%
\indexS{calor bullae}
id quaedam quoque effervescentia%
\indexS{effervescentia}
vel fermentatio%
\indexS{fermentatio}
vel alia causa adventitia%
\indexS{causa adventitia}
aerem producens%
\indexS{causa producens aerem}
ac dilatans%
\indexS{causa dilatans aerem}%
\lbrack,\rbrack\
efficere posset.%
}{%
\lemma{frigidum}%
\Bfootnote{%
\bars{nescio quo causu \textit{erg.}} transvectam %
\textit{(1)}~ponamus, %
\textit{(2)}~, calore suo %
\textit{(a)}~aliquandiu servato; %
\textit{(b)}~conservato; %
\textit{(aa)}~aut nisi %
\textit{(bb)}~an %
\textit{(aaa)}~putemus ipsu %
\textit{(bbb)}~ponamus %
\textit{(ccc)}~ponemus initio %
\lbrack...\rbrack\ externum et %
\textit{(aaaa)}~inclusum %
\textit{(bbbb)}~aerem bullae %
\lbrack...\rbrack\ calorem conservasse. %
\bars{%
\textit{(1)}~An dicemus aerem bullae inclusum %
diutius calorem %
passum antequam dilataretur, ob\indexS{aer dilatatus}\indexS{aer passus calorem} %
\textit{(a)}~pelliculae dil\textlangle a\textrangle\ %
\textit{(b)}~bullae\lbrack\textit{!}\rbrack\ ambientis renisum;\indexS{renisus} %
\textit{(aa)}~aut %
\textit{(bb)}~cum aer ambiens a solo pondere\indexS{pondus aeris superstantis} %
\textit{(aaa)}~incumbentis %
\textit{(bbb)}~superstantis (quod etiam bullae incumbentis) dilatationis impedimentum sentiret\indexS{dilatatio aeris}\indexS{aer superstans}\indexS{aer incumbens bullae} %
\textit{(2)}~An potius %
\lbrack...\rbrack\ intra aquam %
\textit{(a)}~natum %
\textit{(b)}~produci uti %
\textit{(aa)}~aquam ex %
\textit{(bb)}~aerem ex %
\lbrack...\rbrack\ aerem autem %
\bars{sic \textit{erg.}} natum %
\textit{(aaa)}~longe antiquo %
\textit{(bbb)}~et in bulla %
\lbrack...\rbrack\ deest, vi dilatativa %
\textit{(aaaa)}~caloris %
\textit{(bbbb)}~vel explosiva %
\lbrack...\rbrack\ comprimatur. Porro %
\textit{erg.}%
} %
quod calor efficit, %
\textit{(aaaaa)}~idem inter %
\textit{(bbbbb)}~id quaedam %
\textit{(aaaaa-a)}~fermentatio seu explosio %
\textit{(bbbbb-b)}~\bars{quoque \textit{erg.}} effervescentia vel fermentatio %
\textit{(aaaaa-aa)}~efficere posset %
\textit{(bbbbb-bb)}~vel %
\textit{(aaaaa-aaa)}~alias %
\textit{(bbbbb-bbb)}~alia causa %
\textit{(aaaaa-aaaa)}~aerem dilatans efficere posset. %
\textit{(bbbbb-bbbb)}~adventitia aerem %
\lbrack...\rbrack\ efficere posset.%
~\textit{L}%
}}%
\pend%
\pstart%
\textls{Secundo}
\edtext{%
posita majore aeris intra bullam raritate%
\indexS{raritas aeris intra bullam}%
\indexS{bulla vaporis}
ejusque resistentia%
\indexS{resistentia bullae ad compressionem}%
}{%
\lemma{posita}%
\Bfootnote{%
\textit{(1)}~\bars{jam \textit{gestr.}} resistentia bullae aerem rarum continentis %
\textit{(2)}~majore aeris %
\lbrack...\rbrack\ ejusque resistentia%
~\textit{L}%
}}
contra compressionem aeris externi,%
\indexS{compressio bullae ab aere externo}%
\indexS{aer externus}
adeoque posito
\edtext{%
dari et durare bullas%
\indexS{bulla vaporis}
quibus inclusus%
}{%
\lemma{dari}%
\Bfootnote{%
\bars{et durare \textit{erg.}} bullas %
\bars{in \textit{gestr.}} quibus inclusus%
~\textit{L}%
}}
aer sit rarior aere ambiente;%
\indexS{aer inclusus bullae}%
\indexS{aer rarior ambiente}
non putem a vi elastica%
\indexS{vis elastica aeris ambientis}%
\indexS{aer ambiens}
\edtext{%
aeris ambientis%
}{%
\lemma{aeris ambientis}%
\Bfootnote{%
\textit{erg.\;L}%
}}
ascensionem talium bullarum peti%
\indexS{ascensio bullae}
\edtext{%
posse:
tum quia vis Elastica%
\indexS{vis elastica aeris ambientis}
tam sursum quam deorsum
aut ad latus agit;%
}{%
\lemma{posse}%
\Bfootnote{%
\textit{(1)}~, quia vis elastica %
\textit{(2)}~: tum quia vis Elastica %
\bars{non \textit{gestr.}} tam sursum quam deorsum 
\bars{aut ad latus \textit{erg.}} agit;%
~\textit{L}%
}}
tum etiam quia revera tanta est
\edtext{%
vis dilatativa%
\indexS{vis dilatativa aeris}%
}{%
\lemma{vis}%
\Bfootnote{%
\textit{(1)}~Elastica %
\textit{(2)}~dilatativa%
~\textit{L}%
}}
intra bulla%
\indexS{bulla vaporis}
(pro ipsius magnitudine)%
\indexS{magnitudo aeris}
quanta in aere ambiente,%
\indexS{aer ambiens}
eo
\edtext{%
ipso
quia hujus compressioni bulla resistit;%
\indexS{bulla resistens compressioni}%
\indexS{compressio bullae ab aere externo}%
}{%
\lemma{ipso quia}%
\Bfootnote{%
\textit{(1)}~ejus pressioni resistit %
\textit{(2)}~hujus %
\textit{(3)}~hujus compressioni bulla resistit;%
~\textit{L}%
}}
quia nempe
quod vi elasticae aeris%
\indexS{vis elastica aeris rari}
\edtext{%
rari,
id est vi dilatatrici insitae,%
\indexS{vis dilatatrix insita aeri}%
\indexS{vis insita}
defit,%
}{%
\lemma{rari}%
\Bfootnote{%
\textit{(1)}~defit %
\textit{(2)}~, id est %
\lbrack...\rbrack\ insitae, defit,%
~\textit{L}%
}}
id calor supplet%
\indexS{calor supplens vim}
\edtext{%
vel alia quaecunque vis dilatandi superveniens.%
\indexS{vis dilatandi superveniens aeri}%
\indexS{vis superveniens}
Itaque omnino ascensus bullae talis%
\indexS{ascensus bullae}%
\indexS{bulla ascendens in aere}
a pondere%
\indexS{pondus bullae}%
}{%
\lemma{vel}%
\Bfootnote{%
\textit{(1)}~quaecunque %
\textit{(2)}~qualis dixi effervescentia, vim delata %
\textit{(3)}~alia quaecunque vis %
\bars{se \textit{gestr.}} dilatandi superveniens. Itaque %
\textit{(a)}~ascensio %
\textit{(b)}~omnino %
\textit{(aa)}~ascendentis pon %
\textit{(bb)}~ascensus bullae talis a pondere%
~\textit{L}%
}}
petendus erit,
cum nempe tota bulla%
\indexS{bulla vaporis}
composita ex
\edtext{%
pellicula aquea%
\indexS{pellicula aquea bullarume}
\ledpb%
et aere rariore incluso%
\indexS{aer inclusus bullae}%
\indexS{aer rarior ambiente}
minus ponderat quam aer ambiens%
\indexS{aer ambiens}
minus rarus%
\indexS{bulla ponderans minus aere ambiente}
paris cum bulla spatii.%
}{%
\lemma{pellicula}%
\Bfootnote{%
\textit{(1)}~rariore incluso %
\textit{(2)}~aquea %
\textit{(a)}~ex a %
\textit{(b)}~et aere rariore incluso %
\textit{(aa)}~plus %
\textit{(bb)}~minus ponderat quam aer ambiens %
\textit{(aaa)}~paris cum bulla spatii %
\textit{(bbb)}~minus rarus %
\textit{(aaaa)}~paris cum bulla spatii. %
\textit{(bbbb)}~paris cum bulla spatii.%
~\textit{L}%
}}
\pend%
\pstart%
\edtext{%
Itaque
ut bullae ex aqua%
\indexS{bulla ex aqua}
aerem includente%
\indexS{aqua}%
\indexS{bulla includens aerem}%
\indexS{aer inclusus bullae}
vim in aere ambiente ascendendi%
\indexS{vis ascendendi in aere}%
\indexS{aer ambiens}
habeant%
}{%
\lemma{Itaque}%
\Bfootnote{%
\textit{(1)}~si volumus bullas %
\textit{(a)}~ex aqua %
\textit{(b)}~e %
\textit{(c)}~aquae aerem %
\textit{(d)}~includentes %
\textit{(2)}~ut bullae ex aqua aerem %
\textit{(a)}~includente in aere attollantur %
\textit{(b)}~inte %
\textit{(c)}~includente vim %
\lbrack...\rbrack\ ascendendi habeant%
~\textit{L}%
}}%
\lbrack,\rbrack\
oportet
\edtext{%
ut calorem%
\indexS{calor aeris inclusi}
vel aliam vim dilatatricem extraneam,%
\indexS{vis dilatatrix insita aeri}%
\indexS{vis extranea}
calori analogam,%
\indexS{vis analoga caloris}
aeris inclusi raritatem%
\indexS{raritas aeris inclusi}
contra compressionem aeris externi%
\indexS{compressio aeris externi}
sustentare statuamus,%
}{%
\lemma{ut}%
\Bfootnote{%
\textit{(1)}~a calore vel aliqua effervescentia\indexS{effervescentia} %
\textit{(2)}~calorem vel %
\textit{(a)}~aliquam effervescentiam aut analogum aliquod petamus %
\textit{(b)}~aliam vim %
\lbrack...\rbrack\ aeris inclusi %
\textit{(aa)}~rarioris %
\textit{(bb)}~raritatem contra %
\lbrack...\rbrack\ sustentare statuamus,%
~\textit{L}%
}}
ascensum autem a pondere%
\indexS{ascensus bullae}%
\indexS{pondus bullae}
\edtext{%
derivemus,
sed pelliculas%
\indexS{pellicula aquea bullarume}%
}{%
\lemma{derivemus}%
\Bfootnote{%
\textit{(1)}~, pelliculas autem %
\textit{(2)}~. Porro pelliculas %
\textit{(3)}~, sed pelliculas%
~\textit{L}%
}}
admodum tenues esse
necesse erit%
\lbrack;\rbrack\
quod ita
\edlabel{LH_37_03_123-124_123r_solmu-1}%
determinabimus:%
\edtext{%
}{%
\xxref{LH_37_03_123-124_123r_solmu-1}{LH_37_03_123-124_123r_solmu-2}%
\lemma{determinabimus:}%
\Bfootnote{%
\textit{(1)}~Sit %
\textit{(2)}~Habeat%
~\textit{L}%
}}
\edtext{}{%
\lemma{\textit{Am Satzende:}}%
\Afootnote{%
Pagina\indexS{pagina}%
\textsuperscript{\lbrack a\rbrack}
quarta abhinc signo\indexS{signum} \hebr{'}%
\textsuperscript{\lbrack b\rbrack}%
\\[1.5ex]%
{\footnotesize%
\textsuperscript{\lbrack a\rbrack}~%
Pagina %
\textit{(1)}~3tia %
\textit{(2)}~quarta%
~\textit{L}
\quad%
\textsuperscript{\lbrack b\rbrack}~%
Pagina \lbrack...\rbrack\ signo \hebr{'}:
Siehe
Bl.\:\href{\Scan{LH_37_03_124v}}{124\,v\textsuperscript{o}}\! %
(\refpassage{LH_37_03_123-124_124v_viertestufe-1}{LH_37_03_123-124_124v_viertestufe-2}).
Die gestr. Variante \textit{Pagina 3tia} verweist auf Bl.\:\href{\Scan{LH_37_03_124r}}{124\,r\textsuperscript{o}}\! %
(\refpassage{LH_37_03_123-124_124r_drittestufe-1}{LH_37_03_123-124_124r_drittestufe-2}).%
}%
}}
\pend%
\begin{edfigure}[0.75em][1.50em]%
\includegraphics%
{gesamttex/edit_VIII-4/images/LH_37_03_123-124_d_123r.pdf}
\edcaption[0.25em]{%
\edtext{%
\lbrack\textit{Fig.\:1}\rbrack%
}{%
\lemma{\lbrack\textit{Fig.\:1}\rbrack}%
\Cfootnote[nonum]{%
In \textit{L} (Bl.\:\href{\Scan{LH_37_03_123r}}{123\,r\textsuperscript{o}}) ist ferner ein nicht wiedergegebener Entwurf zum Diagramm überliefert.%
}}%
}%
\pagelabel{LH_37_03_123-124_123r_Fig1}%
\end{edfigure}
\pstart%
\noindent%
\lbrack%
\textit{Erster Ansatz:}%
\rbrack\
\edlabel{LH_37_03_123-124_123rv_erstestufe-1}%
Habeat%
\edlabel{LH_37_03_123-124_123r_solmu-2}
aeris%
\indexS{aer inclusus bullae}
\edtext{%
inclusi sphaera%
\indexS{sphaera}
radium \textit{AB},
tota bulla radium \textit{AC}.%
\indexS{radius bullae}%
\indexS{bulla vaporis}%
}{%
\lemma{inclusi}%
\Bfootnote{%
\textit{(1)}~diameter %
\textit{(2)}~sphaera %
\textit{(a)}~diametro %
\textit{(b)}~radium \textit{AB}, tota bulla %
\textit{(aa)}~diametri %
\textit{(bb)}~radium \textit{AC}.%
~\textit{L}%
}}
Ratio densitatis %
in aere ambiente%
\indexS{densitas aeris ambientis}
ad densitatem %
aeris in bulla sit%
\indexS{densitas aeris in bulla}
ut \textit{d}
ad 1. %
Pondus aquae%
\indexS{pondus aquae}
ad aerem%
\indexS{ratio ponderis aquae ad aerem}
ut 800 ad 1.
Pondus aeris in \textit{AB} sit \textit{p}.
Pelliculae spatium%
\indexS{spatium pelliculae aqueae}%
\indexS{pellicula aquea bullarume}
ad aeris inclusi spatium%
\indexS{spatium aeris inclusi}%
\indexS{aer inclusus bullae}
ut $AC\,\text{cub.} - AB\,\text{cub.}$ ad \textit{AB}\,cub.
Pelliculae pondus%
\indexS{pondus pelliculae aqueae}%
\indexS{pellicula aquea bullarume}
ad aeris inclusi%
\indexS{pondus aeris in bulla}%
\indexS{aer inclusus bullae}
\edtext{%
pondus \textit{p}
in ratione composita ex%
}{%
\lemma{pondus \textit{p}}%
\Bfootnote{%
\textit{(1)}~ut %
\textit{(2)}~in ratione composita ex%
~\textit{L}%
}}
hac et ratione 800 ad 1:
adeoque habebis
\bracks{%
\href{\Scan{LH_37_03_123v}}%
{123\,v\textsuperscript{o}}%
}
pelliculae%
\indexS{pondus pelliculae aqueae}%
\indexS{pellicula aquea bullarume}
\edtext{%
pondus,
si 800\textit{p} multiplices
per $AC\,\text{cub.} - AB\,\text{cub.}$%
}{%
\lemma{pondus,}%
\Bfootnote{%
\textit{(1)}~$p(AC\,\text{cub.} - AB\,\text{cub.})$ %
\textit{(2)}~\textit{p} multiplicat. %
\textit{(3)}~si \textit{p} %
\textit{(4)}~si 800\textit{p} multiplices %
\textit{(a)}~per $AC - A$ %
\textit{(b)}~per $AC\,\text{cub.} - AB\,\text{cub.}$%
~\textit{L}%
}}
et
\edtext{%
proven\lbrack i\rbrack ens%
}{%
\lemma{proventens\lbrack\textit{!}\rbrack}%
\Bfootnote{%
\textit{L\;ändert Hrsg.}%
}%
\lemma{proven\lbrack i\rbrack ens}%
\Cfootnote{%
Vgl. zur Berichtigung \refpassage{LH_37_03_123-124_123r_proveniens-1}{LH_37_03_123-124_123r_proveniens-2}.%
}%
}
dividas
\edtext{%
per \textit{AB}\,cub.%
\edtext{}{\lemma{\textit{Am Rand:}}\Afootnote{%
$\displaystyle\frac{800p, \, AC\,\text{cub.} - AB\,\text{cub.}}{AB\,\text{cub.}} + \frac{AC\,\text{cub.}}{AB\,\text{cub.}}p$%
}}
Quod si ad pondus hoc pelliculae%
\indexS{pondus pelliculae aqueae}%
\indexS{pellicula aquea bullarume}
addatur \textit{p} pondus aeris inclusi%
\indexS{pondus aeris in bulla}%
\indexS{aer inclusus bullae}
prodibit pondus bullae totius%
\indexS{pondus bullae}%
\indexS{bulla vaporis}%
\lbrack,\rbrack\
quod etiam prodibit %
\edlabel{LH_37_03_123-124_123v_799p-1}%
si a 800\textit{p}
multiplicato per \textit{AC}\,cub.
et diviso per \textit{AB}\,cub.
detrahas 799\textit{p}.%
\edlabel{LH_37_03_123-124_123v_799p-2}%
}{%
\lemma{per \textit{AB}\,cub.}%
\Bfootnote{%
\textit{(1)}~Sed totam  %
\textit{(2)}~Cui %
\textit{(3)}~Quod si %
\lbrack...\rbrack\ aeris inclusi %
\textit{(a)}~quod e %
\textit{(b)}~et %
\textit{(c)}~prodibit pondus bullae totius %
\textit{(aa)}~8 %
\textit{(bb)}~80 %
\textit{(cc)}~799\textit{p} %
\textit{(dd)}~800\textlangle\textendash\textrangle\ %
\textit{(ee)}~quod e %
\textit{(ff)}~quod etiam prodibit %
\textit{(aaa)}~si a 799\textit{p} detrahas id quod %
\textit{(aaaa)}~prodit %
\textit{(bbbb)}~fit 800\textit{p} multiplicando per \textit{AB}\,cub. et productum deinde per \textit{AC} cub. %
\textit{(bbb)}~si a 800\textit{p} %
\lbrack...\rbrack\ detrahas 799\textit{p}.%
~\textit{L}%
}}
At
\edtext{%
aer ambiens%
\indexS{aer ambiens}
paris cum bulla spatii%
\indexS{spatium bullae}%
}{%
\lemma{aer}%
\Bfootnote{%
\textit{(1)}~paris spatii, s %
\textit{(2)}~ambien %
\textit{(3)}~ambiens paris cum bulla spatii%
~\textit{L}%
}}
erit ad pondus aeris in bulla%
\indexS{pondus aeris in bulla}%
\indexS{aer inclusus bullae}%
\indexS{bulla vaporis}
\edtext{%
inclusi
in ratione composita
ex ratione densitatium%
\indexS{densitas aeris}
quae est \textit{d} ad 1,
et ex ratione spatiorum,
nempe spatii bullae totius%
\indexS{spatium bullae}
ad spatium aeris inclusi%
\indexS{spatium aeris inclusi}
seu \textit{AC}\,cub.
ad cub.\,\textit{AB}.%
}{%
\lemma{inclusi in}%
\Bfootnote{%
\textit{(1)}~composita %
\textit{(2)}~ratione composita %
\textit{(a)}~\textit{d} ad 1, et spatii bullae totius ad spatium aeris inclusi, seu cubi \textit{AC} ad cubum \textit{AB} %
\textit{(b)}~ex ratione %
\lbrack...\rbrack\ ad 1, et %
\bars{ex \textit{erg.}} ratione spatiorum, %
\lbrack...\rbrack\ aeris inclusi seu %
\textit{(aa)}~cub.\,\textit{AB} ad %
\textit{(bb)}~\textit{AC}\,cub. ad cub.\,\textit{AB}.%
~\textit{L}%
}}
Ergo aeris ambientis%
\indexS{pondus aeris ambientis}%
\indexS{aer ambiens}
\edtext{%
pari\lbrack s\rbrack%
}{%
\lemma{pari}%
\Bfootnote{%
\textit{L\;ändert Hrsg.}%
}}
cum bulla spatii pondus%
\indexS{spatium bullae}%
\indexS{bulla vaporis}
prodit,
\edtext{%
si productum ex \textit{p} in \textit{d}
multiplices porro per \textit{AC}\,cub.%
}{%
\lemma{si}%
\Bfootnote{%
\textit{(1)}~\textit{dp} in %
\textit{(2)}~productum ex %
\lbrack...\rbrack\ per \textit{AC}\,cub.%
~\textit{L}%
}}
et
\edlabel{LH_37_03_123-124_123r_proveniens-1}
proveniens
\edlabel{LH_37_03_123-124_123r_proveniens-2}%
dividas per \textit{AB}\,cub.
\edlabel{LH_37_03_123-124_123v_randrechnung-1}%
\edtext{}{%
\xxref{LH_37_03_123-124_123v_randrechnung-1}{LH_37_03_123-124_123v_randrechnung-2}%
\lemma{\textit{Am Rand:}}%
\Afootnote{%
\rule[0mm]{0mm}{12mm}%
$\displaystyle\frac{pd \cdot AC\,\text{cub.}}{AB\,\text{cub.}} \groesser \displaystyle\frac{800p, \, AC\,\text{cub.} - AB\,\text{cub.}}{AB\,\text{cub.}}$
\quad \ %
$\displaystyle d - 800, \, \frac{AC\,\text{cub.}}{AB\,\text{cub.}} \groesser 1$%
}}%
Itaque
hoc pondus%
\indexS{pondus aeris ambientis}
debet esse majus
quam pondus bullae,%
\indexS{pondus bullae}
ut bulla in aere ascendat,%
\indexS{bulla ascendens in aere}%
\indexS{bulla vaporis}
\ledpb%
seu
$\displaystyle\frac{dp \cdot AC\,\text{cub.}}{AB\,\text{cub.}}$
maj. quam
$\displaystyle\frac{800p \cdot AC\,\text{cub.}}{AB\,\text{cub.}} - 
\edtext{%
\edtext{%
800p$.%
\edlabel{LH_37_03_123-124_123v_randrechnung-2}%
}{%
\lemma{$-\,800p$}%
\Cfootnote{%
Der Subtrahend ist eigentlich 799\textit{p}
(vgl. \refpassage{LH_37_03_123-124_123v_799p-1}{LH_37_03_123-124_123v_799p-2}).
Der Fehler wirkt sich auf die folgenden Ableitungen bis zu \refpassage{LH_37_03_123-124_123rv_erstestufe-2}{LH_37_03_123-124_123rv_erstestufe-2} 
aus.
Weitere Fehler kommen hinzu.%
}}
Adeoque 800,
vel quisquis sit numerus,
debet esse major quam numerus%
\rule[0mm]{0mm}{6mm}
$\displaystyle\frac{800 - d, \, AC\,\text{cub.}}{AB\,\text{cub.}}$,
ubi 800 est numerus%
\indexS{numerus exprimens rationem}
exprimens rationem gravitatis specificae aquae%
\indexS{gravitas specifica aquae}%
\indexS{aqua}%
\rule[0mm]{0mm}{4mm}
ad gravitatem specificam aeris ordinarii,%
\indexS{gravitas specifica aeris}%
\indexS{aer ordinarius}
et \textit{d} est numerus exprimens rationem%
\indexS{numerus exprimens rationem}
gravitatis specificae seu densitatis aeris ordinarii%
\indexS{densitas aeris ordinarii}
ad densitatem aeris dilatati%
\indexS{aer dilatatus}
in bulla inclusi.%
\indexS{densitas aeris in bulla}
Itaque debet ratio 800 ad $800 - d$ major esse
quam ratio \textit{AC}\,cub. ad \textit{AB}\,cub.
\edlabel{LH_37_03_123-124_123v_nbnrchng-1}%
\edtext{}{%
\xxref{LH_37_03_123-124_123v_nbnrchng-1}{LH_37_03_123-124_123v_nbnrchng-2}%
\lemma{Ponamus \lbrack...\rbrack\ $\sqrt[3]{\vphantom{()}}(8000 - 100)$}%
\Cfootnote{%
Vgl. hierzu die auf S.\:\pageref{LH_37_03_123-124_123v_nbnrchngn}\,f. abgedruckten Nebenrechnungen.}}%
Ponamus \textit{d} esse 10
seu aerem ambientem%
\indexS{aer ambiens}
esse decuplo graviorem incluso,%
\indexS{aer inclusus bullae}
debet major esse ratio 20
ad radicem cubicam ex 790%
\indexS{radix cubica}
quam ratio radii \textit{AC}
ad radium \textit{AB};%
\indexS{radius bullae}
adeoque eo casu crassities pelliculae $AC - AB$%
\indexS{crassities pelliculae}
erit ad radium sphaerae aeris inclusi%
\indexS{sphaera aeris inclusi}%
\indexS{aer inclusus bullae}
minor quam ratio
$\displaystyle20 - \sqrt[3]{\vphantom{j}}(8000 - 100)$%
\edlabel{LH_37_03_123-124_123rv_erstestufe-2}%
\edlabel{LH_37_03_123-124_123v_nbnrchng-2}%
}{%
\lemma{800\textit{p}.}%
\Bfootnote{%
\textit{(1)}~Seu 800 %
\textit{(2)}~Adeoque 800, %
\textit{(a)}~ratio po %
\textit{(b)}~numerus %
\textit{(c)}~vel quisquis sit numerus %
\textit{(aa)}~exprimens rationem %
\textit{(aaa)}~ponderis specifici aquae ad pondus specific\indexS{pondus specificum aquae} %
\textit{(bbb)}~gravitatis specificae aquae et aeris, %
\textit{(bb)}~, debet esse %
\lbrack...\rbrack\ numerus \rule[0mm]{0mm}{4mm} $\displaystyle\frac{800 - d, \, AC\,\text{cub.}}{AB\,\text{cub.}}$, %
\textit{(aaa)}~vel ratio 800 ad 1 debet esse major quam ratio composita ex ratione $800 - 1$ ad 1 et ratione \textit{AC}\,cub. ad \textit{AB}\,cub., posito %
\textit{(aaaa)}~\textit{d} esse rat %
\textit{(bbbb)}~rationem \textit{d} ad 1 esse rationem %
\textit{(aaaaa)}~aeris %
\textit{(bbbbb)}~densitatis seu gravitatis specificae aeris ambientis %
\textit{(aaaaa-a)}~ad inclusi densi %
\textit{(bbbbb-b)}~et inclusi %
\textit{(aaaaa-aa)}~. Si duo numeri essent aequales, foret %
\textit{(bbbbb-bb)}~. Si essent aequa %
\textit{(ccccc-cc)}~%
, adeoque ratio 800 ad $800 - d$
\textit{(bbb)}~ubi \textit{d} est numerus %
\textit{(ccc)}~ubi 800 %
\lbrack...\rbrack\ specificae aquae %
\textit{(aaaa)}~ad gr. spec. aeris %
\textit{(bbbb)}~ad gravitatem specificam aeris ordinarii, %
\textit{(aaaaa)}~et \textit{d} est est n %
\textit{(bbbbb)}~et \textit{d} est numerus exprimens %
\textit{(aaaaa-a)}~densitatem in qua %
\textit{(bbbbb-b)}~rationem gravitatis %
\lbrack...\rbrack\ graviorem incluso, %
\textit{(aaaaa-aa)}~fiet %
\textit{(bbbbb-bb)}~debet ratio bullae %
\textit{(ccccc-cc)}~20 ad radicem
\textit{(ddddd-dd)}~debet major esse %
\textit{(aaaaa-aaa)}~numer %
\textit{(bbbbb-bbb)}~ratio 20 %
\lbrack...\rbrack\ ex 790 %
\textit{(aaaaa-aaaa)}~quam est $\displaystyle\sqrt[3]{\vphantom{j}}790$
\textit{(bbbbb-bbbb)}~quam %
\textit{(aaaaa-aaaaa)}~\textit{AC} ad %
\textit{(bbbbb-bbbbb)}~ratio radii \textit{AC} ad %
\textit{(aaaaa-aaaaa-a)}~rationem %
\textit{(bbbbb-bbbbb-b)}~radium \textit{AB}; %
\lbrack...\rbrack\ aeris inclusi %
\textit{(aaaaa-aaaaa-aa)}~ut 800 %
\textit{(bbbbb-bbbbb-bb)}~minor quam ratio $\displaystyle20 - \sqrt[3]{\vphantom{j}}(8000 - 100)$%
~\textit{L}%
}}
\lbrack%
\textit{der Text bricht ab und wird durch eine Trennlinie abgeschlossen,
die den gesamten Abschnitt \refpassage{LH_37_03_123-124_123rv_erstestufe-1}{LH_37_03_123-124_123rv_erstestufe-2} umkreist.}%
\rbrack%
\pend%
\pstart%
Hinc etiam intelligi%
\indexS{causa intellecta}
\edtext{%
causa potest cur%
}{%
\lemma{causa}%
\Bfootnote{%
\textit{(1)}~est c %
\textit{(2)}~potest cur%
~\textit{L}%
}}
bullae rursus descendant,%
\indexS{descensus bullae}%
\indexS{bulla descendens}%
\indexS{bulla vaporis}
nam cessante vi dilatatrice%
\indexS{vis dilatatrix adventitia aeri}%
\indexS{vis adventitia}
adventitia %
comprimentur a vi elastica aeris%
\indexS{bulla compressa a vi elastica aeris}%
\indexS{vis elastica aeris}
\edtext{%
ambientis,%
\indexS{aer ambiens}
adeoque imminuto volumine%
\indexS{volumen bullae imminutum}
descendere nitentur,%
}{%
\lemma{ambientis,}%
\Bfootnote{%
\textit{(1)}~adeoque ubi aer inclusus et ambiens paris densitatis erunt, %
\textit{(2)}~adeoque in eo descendent, %
\textit{(3)}~adeoque imminuto volumine descendere nitentur,%
~\textit{L}%
}}
vel dirumpentur,%
\indexS{bulla dirupta}
quo casu
\edtext{%
etiam aqua per se sumta%
\indexS{aqua descendens in aere}
utique in aere descendere nitetur.%
}{%
\lemma{etiam}%
\Bfootnote{%
\textit{(1)}~descendet aqua %
\textit{(a)}~ad %
\textit{(b)}~de %
\textit{(c)}~per %
\textit{(2)}~aqua per se sumta utique %
\textit{(a)}~in guttas %
\textit{(b)}~in aere descendet. %
\textit{(c)}~in aere descendere nitetur.%
~\textit{L}%
}}
\ledpb%
Quod si tantae sint%
\indexS{tenuitas bullulae}
\edtext{%
tenuitatis %
vel bullulae ex aere et aqua%
\indexS{bullula ex aere et aqua}%
}{%
\lemma{tenuitatis}%
\Bfootnote{%
\textit{(1)}~guttae vel laminae aqu %
\textit{(2)}~vel bullulae %
\lbrack...\rbrack\ et aqua%
~\textit{L}%
}}
vel particulae aqueae,%
\indexS{tenuitas particulae aqueae}%
\indexS{particulae aqueae}
\edtext{%
ut non facile aerem%
\indexS{aer non carens cohaesione}
cohaesione aliqua partium suarum non carentem%
\indexS{cohaesio partium aeris}
penetrent,
adeoque in eo%
\lbrack,\rbrack\ licet graviores%
\lbrack,\rbrack\
natent%
\indexS{bullula natans in aere}
et huc illuc motu ejus jactentur,%
\indexS{motus aeris}
quemadmodum et pulvisculis fieri%
\indexS{pulvisculus}
videmus in aere natantibus;
hoc posito
tamdiu poterit suspendi descensus,%
\indexS{descensus guttae in aere}%
}{%
\lemma{ut}%
\Bfootnote{%
\textit{(1)}~qualecunque %
\textit{(2)}~descendant, quemadmodum et de pulvisculis fieri videmus, sed facile huc illuc motu ejus jactentur, %
\textit{(a)}~si %
\textit{(b)}~fiet ta 
\textit{(c)}~alio %
\textit{(3)}~non facile %
\textit{(a)}~in aere descendant, sed %
\textit{(b)}~aerem %
\textit{(aa)}~penetrent %
\textit{(bb)}~cohaesione aliqua %
\lbrack...\rbrack\ poterit suspendi %
\bars{in motu \textit{erg. u. gestr.}} descensus,%
~\textit{L}%
}}
donec plures guttae%
\indexS{guttae concrescentes}%
\indexS{gutta descendens in aere}
\edtext{%
minores concrescant in majores,
quae resistentiam medii vincant;%
\indexS{resistentia medii}%
}{%
\lemma{minores}%
\Bfootnote{%
\textit{(1)}~in %
\textit{(2)}~concrescant in majores, %
\textit{(a)}~ubi %
\textit{(b)}~quo facto %
\textit{(c)}~quae resistentiam medii vincant;%
~\textit{L}%
}}
quoniam resistentiae%
\indexS{resistentia medii}
crescunt tantum
ut superficies guttarum,%
\indexS{superficies guttae}
sed pondera crescent ut earum%
\indexS{pondus guttae}
\edtext{%
soliditates,%
\indexS{soliditas guttae}
sunt autem%
}{%
\lemma{soliditates,}%
\Bfootnote{%
\textit{(1)}~nam %
\textit{(2)}~sunt autem%
~\textit{L}%
}}
\edtext{%
superficies ut quadrata,
at soliditates ut cubi diametrorum.%
\indexS{diameter guttae}%
}{%
\lemma{superficies \lbrack...\rbrack\ diametrorum}%
\Cfootnote{%
Siehe \textsc{Euklid}, \textit{Elementa}, lib.\;XII, prop.\;18.\cite{01051}%
}}
\bracks{%
\href{\Scan{LH_37_03_123r}}%
{123\,r\textsuperscript{o}}%
}%
\pend%
\vspace{1.750em}%
\footnotesize%
\pstart%
\noindent%
\lbrack%
{\normalsize%
\textit{Zweiter Ansatz, kleingedruckter Text gestrichen:}%
\rbrack%
}
\edlabel{LH_37_03_123-124_123r+124r_zweitestufe-1}%
Ponamus %
800 esse numerum
quo aqua%
\indexS{aqua gravior aere}
est gravior aere ordinario,%
\indexS{aer ordinarius}
et \textit{d} esse numerum
quo aer ordinarius seu ambiens%
\indexS{aer ambiens}
est gravior
aere intra pelliculam incluso%
\indexS{aer inclusus bullae}%
\indexS{pellicula aquea bullarume}%
\lbrack;\rbrack\
dico
crassitiem pelliculae aqueae%
\indexS{crassities pelliculae aqueae}%
\indexS{pellicula aquea bullarume}
se haber\textlangle e\textrangle\
ad radium sphaerae aeris%
\indexS{sphaera aeris inclusi}%
\indexS{aer inclusus bullae}
\edtext{%
inlcusi %
in ratione minore quam est
$\displaystyle 20 - \sqrt[3]{\vphantom{j}}(8000 - \langle100)\rangle$%
}{%
\lemma{inlcusi}%
\Bfootnote{%
\textit{(1)}~ut %
\textit{(a)}~8000 %
\textit{(b)}~$\displaystyle 20 - \sqrt[3]{\vphantom{j}}(8000 - \langle100)\rangle$ se habet %
\textit{(2)}~in ratione minore quam est $\displaystyle 20 - \sqrt[3]{\vphantom{j}}(8000 - \langle100)\rangle$%
~\textit{L}%
}}
ad
$\displaystyle \sqrt[3]{\vphantom{j}}(8000 - 10d).$
Itaque si \textlangle\textit{d}\textrangle\ sit 10,
seu si aer inclusus sit decuplo
\edtext{%
\textlangle rarior\textrangle%
}{%
\lemma{\textlangle rarior\textrangle}%
\Cfootnote{%
Das Wort ist noch tlw. lesbar im S-Film Nr.\:116 der GWLB Hannover.%
}}
ordinario ambiente,%
\indexS{aer rarior ambiente}
erit pelliculae
\edtext{%
aqueae%
}{%
\lemma{aqueae}%
\Bfootnote{%
\textit{erg.\;L}%
}}
crassities%
\indexS{crassities pelliculae aqueae}%
\indexS{pellicula aquea bullarume}
ad radium sphaerae aeris inclusi%
\indexS{sphaera aeris inclusi}%
\indexS{aer inclusus bullae}
in ratione
\edtext{%
minore
quam est
20 $-$ rad.cub. de 7900 ad rad.cub. de 7900,
id est minore quam 1 ad 200.%
}{%
\lemma{minore}%
\Bfootnote{%
\textit{(1)}~est %
\textit{(2)}~quam est %
\textit{(a)}~$\displaystyle 20 - \sqrt[3]{\vphantom{j}}\langle7900\rangle$ ad $\sqrt[3]{\vphantom{j}}7900$ %
\textit{(b)}~20 $-$ rad.cub.de 7900 ad rad.cub.de 7900 %
\textit{(aa)}~id e %
\textit{(bb)}~, id est minore quam 1 ad %
\textit{(aaa)}~40\textlangle0\textrangle\ %
\textit{(bbb)}~200.%
~\textit{L}%
}}
\normalsize%
\bracks{%
\href{\Scan{LH_37_03_124r}}%
{124\,r\textsuperscript{o}}%
}%
\pend%
\vspace{0.50em}%
\normalsize%
\pstart%
Quod sic
\edtext{%
ostendo: %
Centro bullae%
\indexS{centrum bullae}%
\indexS{bulla vaporis}
existente \textit{A},
habeat aeris%
}{%
\lemma{ostendo:}%
\Bfootnote{%
\textit{(1)}~Sit %
\textit{(2)}~Habeat areas\lbrack\textit{!}\rbrack\ %
\textit{(3)}~Centro bullae %
\lbrack...\rbrack\ habeat aeris%
~\textit{L}%
}}
inclusi sphaera%
\indexS{sphaera aeris inclusi}%
\indexS{aer inclusus bullae}
radium \textit{AB},
bulla tota
ex aere et pellicula aquea constans,%
\indexS{bulla constans ex aqua et aere}%
\indexS{pellicula aquea bullarume}
radium
\edtext{%
\textit{AC};
et Gravitas specifica aquae%
\indexS{gravitas specifica aquae}%
\indexS{aqua in bulla}
sit ad Gravitatem specificam aeris%
\indexS{gravitas specifica aeris}
oridinarii seu ambientis%
\indexS{aer ordinarius}%
\indexS{aer ambiens}
ut 800 ad 1,
et gravitas specifica seu densitas aeris ambientis%
\indexS{densitas aeris ambientis}
ad densitatem (multo minorem) aeris inclusi%
\indexS{densitas aeris in bulla}%
\indexS{aer inclusus bullae}
ut \textit{d} ad 1.
Erit gravitas specifica aquae%
\indexS{gravitas specifica aquae}%
\indexS{aqua in bulla}
ad %
gravitatem specificam aeris inclusi%
\indexS{gravitas specifica aeris}%
\indexS{aer inclusus bullae}%
}{%
\lemma{\textit{AC};}%
\Bfootnote{%
\textit{(1)}~centro bullae existe %
\textit{(2)}~sit aeris %
\textit{(3)}~et %
\textit{(4)}~numerus 800 ex %
\textit{(5)}~pondus aeris inclusi sit \textit{p}; pelliculae spatium erit differentia cuborum ab \textit{AC} et ab \textit{AB}. Et pelliculae pondus erit ad \textit{p} in ratione composita\indexS{pondus aeris in bulla}\indexS{spatium pelliculae aqueae}\indexS{pondus pelliculae aqueae} %
\textit{(a)}~aquae ad %
\textit{(aa)}~aerem %
\textit{(bb)}~gravitatem %
\textit{(b)}~aquae et aeris %
\textit{(c)}~spatii aquae ad spatium %
\textit{(6)}~pondus aeris %
\textit{(7)}~pondus aer %
\textit{(8)}~pondus aquae est ad pondus %
\textit{(9)}~gravitas specifica aqueae %
\textit{(a)}~est %
\textit{(b)}~erit %
\textit{(10)}~et Gravitas specifica aquae %
\textit{(a)}~erit ad %
\textit{(b)}~sit ad %
\textit{(c)}~erit %
\textit{(d)}~sit ad Gravitatem specificam aeris %
\textit{(aa)}~ut %
\textit{(bb)}~oridinarii seu %
\lbrack...\rbrack\ aeris inclusi%
~\textit{L}%
}}
in ratione
\edtext{%
composita ex rationibus
aquae ad aerem ordinarium,%
\indexS{aqua in bulla}
\ledpb%
et aeris ordinarii ad inclusum,%
\indexS{aer inclusus bullae}%
\indexS{aer ordinarius}
quae sunt 800 ad 1,
et \textit{d} ad 1,
adeoque erit ut 800\textit{d} ad 1.%
}{%
\lemma{composita}%
\Bfootnote{%
\textit{(1)}~800 ad 1 et \textit{d} ad 1, id est ut 800\textit{d} ad 1, %
\textit{(2)}~ex rationibus %
\textit{(3)}~800 ad 1 %
\textit{(4)}~ex rationibus %
\lbrack...\rbrack\ ut 800\textit{d} ad 1.%
~\textit{L}%
}}
Rursus pondus%
\indexS{pondus pelliculae aqueae}%
\indexS{pellicula aquea bullarume}
\edtext{%
pelliculae aqueae est ad pondus%
}{%
\lemma{pelliculae}%
\Bfootnote{%
\textit{(1)}~est ad pondus %
\textit{(2)}~aqueae est ad pondus%
~\textit{L}%
}}
aeris inclusi%
\indexS{pondus aeris in bulla}
in ratione composita%
\indexS{ratio composita e rationibus}
ex rationibus gravitatis specificae aquae%
\indexS{gravitas specifica aquae}
ad gravitatem specificam aeris inclusi,%
\indexS{gravitas specifica aeris}%
\indexS{aer inclusus bullae}
et spatii aqueae pelliculae%
\indexS{spatium pelliculae aqueae}%
\indexS{pellicula aquea bullarume}
ad spatium aeris%
\indexS{spatium aeris inclusi}%
\indexS{aer inclusus bullae}
\edtext{%
inclusi,
quae rationes sunt%
}{%
\lemma{inclusi,}%
\Bfootnote{%
\textit{(1)}~id est %
\textit{(2)}~quae rationes sunt%
~\textit{L}%
}}
800\textit{d} ad 1,
et
\edtext{%
\textit{AC}\,cubus
demto \textit{AB}\,cubo%
}{%
\lemma{\textit{AC}\,cubus}%
\Bfootnote{%
\textit{(1)}~minus %
\textit{(2)}~demto \textit{AB}\,cubo%
~\textit{L}%
}}
ad \textit{AB}\,cubum,
adeoque ratio erit
ut
\edtext{%
800\textit{d}
numerus multiplicans differentiam cuborum,%
\indexS{differentia cuborum}
divisu\lbrack s\rbrack\
per cubum minorem,%
}{%
\lemma{800\textit{d}}%
\Bfootnote{%
\textit{(1)}~in %
\textit{(2)}~numerus %
\textit{(a)}~dividens %
\textit{(b)}~multiplicans differentiam cuborum, %
\textit{(aa)}~ad cubum %
\textit{(bb)}~\bars{divisum \textit{ändert Hrsg.}} per cubum minorem, %
~\textit{L}%
}}
est ad
\edtext{%
unitatem.%
\indexS{unitas}
Itaque
si pondus aeris inclusi
sit \textit{p},
pondus pelliculae%
\indexS{pondus pelliculae aqueae}%
\indexS{pellicula aquea bullarume}
erit
$800dp,\, (AC\,\text{cub.} - AB\,\text{cub.}) : AB\,\text{cub.}$
Cui si%
}{%
\lemma{unitatem.}%
\Bfootnote{%
\textit{(1)}~Cui si %
\textit{(2)}~Itaque si ad pondus %
\textit{(3)}~Itaque si pondus %
\textit{(a)}~pelliculae sit \textit{p} pondus aeris inclusi erit %
\textit{(b)}~aeris inclusi %
\lbrack...\rbrack\ Cui si%
~\textit{L}%
}}
addatur ipsum pondus aeris inclusi~\textit{p},
prodibit pondus totius bullae%
\indexS{pondus bullae}
ex pellicula et aere compositae.%
\indexS{bulla composita ex aere et pellcula aquea} %
Sed aer
\edtext{%
ordinarius seu ambiens%
\indexS{aer ordinarius}%
\indexS{aer ambiens}%
}{%
\lemma{ordinarius seu ambiens}%
\Bfootnote{%
\textit{erg.\;L}%
}}
paris cum bulla spatii%
\indexS{spatium bullae}
est ad \textit{p}
pondus aeris inclusi%
\indexS{pondus aeris in bulla}%
\indexS{aer inclusus bullae}
\edtext{%
in ratione composita
gravitatum specificarum%
\indexS{gravitas specifica aeris}
et spatiorum,%
\indexS{spatium aeris inclusi}%
}{%
\lemma{in}%
\Bfootnote{%
\textit{(1)}~composita %
\textit{(2)}~ratione composita %
\textit{(a)}~gravitatis specificae %
\textit{(b)}~gravitatum specificarum et spatiorum,%
~\textit{L}%
}}
id est in ratione composita ex rationibus%
\indexS{ratio composita e rationibus}
\textit{d} ad 1 et \textit{AC}\,cub. ad \textit{AB}\,cub.
Ergo pondus aeris ambientis%
\indexS{pondus aeris ambientis}
qui sit paris cum bulla spatii%
\indexS{spatium bullae}
est $dp \cdot AC\,\text{cub.} : AB\,\text{cub.}$ %
Itaque ut bulla in aere ascendere nitatur,%
\indexS{bulla ascendens in aere}
oportet ut pondus bullae sit minus%
\indexS{pondus bullae}
pondere aeris ambientis paris spatii,%
\indexS{pondus aeris ambientis}
seu ut%
\rule[0mm]{0mm}{6mm}
$\displaystyle\frac{800dp \, (AC\,\text{cub.} - AB\,\text{cub.})}{AB\,\text{cub.}} + p$
sit minus quam
$\displaystyle\frac{dp \cdot AC\,\text{cub.}}{AB\,\text{cub.}},$
adeoque ut
$800d \, (AC\,\text{cub.} - AB\,\text{cub.}) + AB\,\text{cub.}$
sit minus quam
\edtext{%
$d \cdot AC\,\text{cub.},$
vel ut
$799d \cdot AC\,\text{cub.}$%
\rule[0mm]{0mm}{4mm}
sit minor quam $(800d - 1) \cdot AB\,\text{cub.},$
vel ut
\textit{AB}\,cubus sit minor quam
$(800\,AB\,\text{cub.} - 799\,AC\,\text{cub.})\,d,$
vel denique
ut \textit{d} sit major%
}{%
\lemma{$d \cdot AC\,\text{cub.}$}%
\Bfootnote{%
\textit{(1)}~vel ut %
\textit{(a)}~800 %
\textit{(b)}~$799d \cdot AC\,\text{cub.}$ sit minus quam $(800d + 1) \cdot AC\,\text{cub.},$ seu denique ut sit %
\textit{(2)}~vel ut %
\lbrack...\rbrack\ $(800d - 1) \cdot AB\,\text{cub.},$ vel ut %
\textit{(a)}~$799d \cdot AC\,\text{cub.}$ %
\textit{(b)}~$(800\,AB\,\text{cub.} - 799\,AC\,\text{cub.})d$ sit minus quam %
\textit{(aa)}~\textit{AB}\,cub. %
\textit{(bb)}~nihilum %
\textit{(c)}~\textit{AB}\,cubus sit %
\lbrack...\rbrack\ ut \textit{d} sit %
\textit{(aa)}~minor %
\textit{(bb)}~major%
~\textit{L}%
}}
quam \textit{AB}\,cub.
divisus per
$800\,AB\,\text{cub.} -
\edtext{%
799\,AC\,\text{cub.}$
Ergo \textit{AC}\,cubi
excessus super \textit{AB}\,cubum
sit pars
$\displaystyle\frac{1}{r}\text{esima}$ cubi\,\textit{AB},
fiet
$\displaystyle AC\,\text{cub.} = AB\,\text{cub.} + \frac{1}{r}AB\,\text{cub.}\;$%
\ledpb%
\lbrack et\rbrack\
fiet
$\displaystyle 800\,AB\,\text{cub.} - 799\,AC\,\text{cub.} = 800\,AB\,\text{cub.} - 799\,AB\,\langle\text{cub.}\rangle - \frac{1}{r}AB\,\text{cub.} = AB\,\text{cub.}\, (1 - \frac{1}{r}).$
Ergo \textit{d} major quam 1 divis. per $\displaystyle1 - \frac{1}{r}$
seu quam \textit{r} divis. per%
\edlabel{LH_37_03_123-124_123r+124r_zweitestufe-2}%
}{%
\lemma{$799\,AC\,\text{cub.}$}%
\Bfootnote{%
\textit{(1)}~Itaque si pellicula\indexS{pellicula aquea bullarume} %
\bars{(exempli gratia) \textit{erg.}} occupet sexagesimam quartam %
\textit{(a)}~spatii partem %
\textit{(b)}~spatii bullae partem, ita ut sit \textit{AC}\,cub. %
\textit{(2)}~Ponamus %
\textit{(a)}~pelliculam esse partem %
\textit{(b)}~spatium pelliculae continere partem \textit{r} spatii bullae, seu esse $\displaystyle\frac{1}{r}AC\,\text{cub.} = AC\,\text{cub.} - AB\,\text{cub.}$, fiet $AB\,\text{cub.} = AC\,\text{cub.}\, (1 - \displaystyle\frac{1}{r}).$ %
\textit{(3)}~Ergo %
\textit{(a)}~sit %
\textit{(aa)}~\textit{d} minor quam ------ seu \textit{d} minor quam %
\textit{(bb)}~\textit{d} major quam 1 divis per (800 - 799r) %
\textit{(b)}~sit $AC\,\text{cub.} = r \cdot AB\,\text{cub.}$ fiet %
\textit{(c)}~\textit{AC}\,cubi excessus %
\lbrack...\rbrack\ fiet $\displaystyle AC\,\text{cub.} = AB\,\text{cub.} + \frac{1}{r}AB\,\text{cub.}$ %
\bars{et \textit{erg. Hrsg.}} fiet $\displaystyle 800\,AB\,\text{cub.} - 799\,AC\,\text{cub.} =$ %
\lbrack...\rbrack\ \textit{r} divis. per%
~\textit{L}%
}}
\lbrack\textit{der Text bricht ab.}\rbrack\
\pend%
\vspace{1.50em}%
\normalsize%
\pstart%
\noindent%
\lbrack%
\textit{Dritter Ansatz:}%
\rbrack\
\edlabel{LH_37_03_123-124_124r_drittestufe-1}%
\edtext{}{%
\xxref{LH_37_03_123-124_124r_drittestufe-1}{LH_37_03_123-124_124r_drittestufe-3}%
\lemma{\hebr{'}}%
\Bfootnote{%
\textit{(1)}~$AB - AC\,\text{cub.} = AB\,\text{cub.} + \displaystyle\frac{1}{r}AB\,\text{cub.},$ erit pelliculae spatium $\displaystyle\frac{1}{r}AB\,\text{cub.},$ et %
\textit{(2)}~Sit %
\textit{(3)}~Gravitas specifica aquae %
\textit{(a)}~erit %
\textit{(b)}~800 vicibus \lbrack...\rbrack\ sphaera aeris inclusi %
\textit{(aa)}~ut cubus %
\textit{(bb)}~radium habeat \textit{AB}, %
\textit{(aaa)}~sphaerae totius bullae radium \textit{AC}, p\textlangle--\textrangle\ erunt %
\textit{(bbb)}~sphaerae totius bullae radium %
\textit{(aaaa)}~\textit{AB} %
\textit{(bbbb)}~erit %
\textit{(cccc)}~\textit{AC}, %
\textit{(aaaaa)}~soliditas %
\textit{(bbbbb)}~\textlangle totum\textrangle\ bullae ad soliditatem co\textlangle--\textrangle\ %
\textit{(ccccc)}~erit spatium \lbrack...\rbrack\ horum cuborum. %
\textit{(aaaaa-a)}~Sit cubus \textit{A} %
\textit{(bbbbb-b)}~Quae sit %
\textit{(aaaaa-aa)}~ut %
\textit{(bbbbb-bb)}~$\displaystyle\frac{1}{r}\,\text{cub.}\,AB,$ erit \text{cub.} %
\bars{\textit{AB} \textit{ändert Hrsg.}} $\displaystyle=\ \text{cub.}\,AB + \frac{1}{r}\,\text{cub.}\,AB.$ %
\textit{(aaaaa-aaa)}~Sit \textit{A} %
\textit{(bbbbb-bbb)}~Jam aeris %
\lbrack...\rbrack\ erit pelliculae %
\bars{aqueae \textit{erg.}} pondus %
\textit{(aaaaa-aaaa)}~$\displaystyle\frac{800d}{rp}\text{cub.}\,AB$ %
\textit{(bbbbb-bbbb)}~ut $\displaystyle 800\, \frac{d}{r}$ %
\bars{\textit{p} \textit{erg. Hrsg.}} cub.\,\textit{AB},%
~\textit{L}%
}}%
\hebr{'}
Gravitas specifica aquae%
\indexS{gravitas specifica aquae}
800 vicibus contineat
gravitatem specificam aeris%
\indexS{gravitas specifica aeris}
ordinarii seu ambientis,%
\indexS{aer ordinarius}%
\indexS{aer ambiens}
et haec \textit{d} vicibus contineat
gravitatem specificam seu densitatem%
\indexS{densitas aeris in bulla}
aeris bullae inclusi.%
\indexS{aer inclusus bullae}%
\indexS{bulla vaporis}
Bullae centrum sit \textit{A},%
\indexS{centrum bullae}
sphaera aeris inclusi%
\indexS{sphaera aeris inclusi}
radium habeat \textit{AB},
sphaerae totius bullae radium \textit{AC},%
\indexS{sphaera bullae}%
\indexS{radius bullae}
erit spatium bullae%
\indexS{spatium bullae}
ad spatium aeris inclusi%
\indexS{spatium aeris inclusi}
ut cubus ab \textit{AC} ad cubum ab \textit{AB},
et spatium pelliculae aqueae%
\indexS{spatium pelliculae aqueae}
ut differentia horum cuborum.%
\indexS{differentia cuborum}
Quae sit
$\displaystyle\frac{1}{r}\,\text{cub.}\,AB,$
erit%
\rule[0mm]{0mm}{5.8mm}
$\displaystyle\text{cub.}\,[AC] = \text{cub.}\,AB + \frac{1}{r}\,\text{cub.}\,AB.$
\edlabel{LH_37_03_123-124_124r_pondusp-1}%
Jam aeris inclusi pondus%
\indexS{pondus aeris in bulla}
sit \textit{p},%
\edlabel{LH_37_03_123-124_124r_pondusp-2}
erit pelliculae aqueae pondus%
\indexS{pondus pelliculae aqueae}%
\indexS{pellicula aquea bullarume}
ut
$\displaystyle 800 \frac{d}{r} [p]\,\text{cub.}\,AB,$%
\edlabel{LH_37_03_123-124_124r_drittestufe-3}
et
\edtext{%
aeris
\edtext{%
inclusi pondus%
\indexS{pondus aeris in bulla}%
\indexS{aer inclusus bullae}
ut
$\displaystyle p \cdot \text{cub.}\,AB,$%
}{%
\lemma{inclusi pondus}%
\Bfootnote{%
\textit{(1)}~\textit{AB} %
\textit{(2)}~ut $\displaystyle p \cdot \text{cub.}\,AB,$%
~\textit{L}%
}}%
}{%
\lemma{aeris \lbrack...\rbrack\ $\displaystyle p \cdot \text{cub.}\,AB$}%
\Cfootnote{%
Das Gewicht der eingeschlossenen Luft ist bereits als \textit{p} definiert (vgl. \refpassage{LH_37_03_123-124_124r_pondusp-1}{LH_37_03_123-124_124r_pondusp-2}).
Die Unstimmigkeit wirkt sich auf die folgenden Ableitungen bis zu \refpassage{LH_37_03_123-124_124r_drittestufe-2}{LH_37_03_123-124_124r_drittestufe-2} aus.%
}}
et bullae pondus%
\indexS{pondus bullae}
ut
$\displaystyle p \cdot \text{cub.}\,AB \cdot (1 + \frac{800d}{r}),$%
\rule[0mm]{0mm}{5.4mm}
et aeris
\edtext{%
ambientis%
}{%
\lemma{ambientis}%
\Bfootnote{%
\textit{erg.\;L}%
}}
paris spatii pondus%
\indexS{pondus aeris ambientis}%
\indexS{aer ambiens}
ut
$\displaystyle 800p \, (1 + \frac{1}{r})\, \text{cub.}\,AB.$
Ergo debet esse
$\displaystyle 800\,(1 + \frac{1}{r})$
majus quam
\edtext{%
$\displaystyle 1 + 800\,\frac{d}{r}$
seu 799 majus quam $\displaystyle 800 \frac{d - 1}{r},$
seu erit ratio
$d - 1$ ad \textit{r}
\ledpb%
minor quam ratio 799 ad 800%
}{%
\lemma{$\displaystyle 1 + 800\,\frac{d}{r}$}%
\Bfootnote{%
\textit{(1)}~seu $\displaystyle 800 + 800\,\frac{1}{r}$ majus quam $\displaystyle 1 + 800\,\frac{d}{r},$ seu %
\textit{(2)}~seu 799 %
\lbrack...\rbrack\ seu erit %
\textit{(a)}~$\displaystyle\frac{799}{8}$ majus quam $d - 1$ divis. per $r - 1$ %
\textit{(b)}~ratio %
\textit{(aa)}~799 ad 800 major quam 
\textit{(aaa)}~ratio %
\textit{(bbb)}~$d - 1$ a %
\textit{(ccc)}~\rule[0mm]{0mm}{4.0mm}ratio %
\textit{(ddd)}~ratio %
\textit{(bb)}~$d - 1$ ad \textit{r} %
\textit{(aaa)}~major %
\textit{(bbb)}~minor quam %
\lbrack...\rbrack\ ad 800%
~\textit{L}%
}}
\edtext{%
\lbrack%
atque%
}{%
\lemma{\lbrack atque}%
\Cfootnote{%
Die eckige Klammer stammt von Leibniz.%
}}
\edtext{%
adeo circiter dici potest rem succedere
si \textit{d} sit minor quam \textit{r},
\rule[0mm]{0mm}{3.4mm}
id est
si \textit{d}
numerus exprimens rationem%
\indexS{numerus exprimens rationem}
densitatis aeris ambientis
ad densitatem aeris inclusi%
\indexS{densitates aeris comparatae}%
\indexS{densitas aeris}
sit minor quam \textit{r}
numerus exprimens rationem%
\indexS{numerus exprimens rationem}
spatii aeris inclusi%
\indexS{spatium aeris inclusi}
ad spatium pelliculae%
\indexS{pellicula aquea bullarume}%
\indexS{spatium pelliculae aqueae}
\edtext{%
aqueae.%
\rbrack%
}{%
\lemma{aqueae.\rbrack\,}%
\Cfootnote{%
Die eckige Klammer stammt von Leibniz.%
}}%
}{%
\lemma{adeo}%
\Bfootnote{%
\textit{(1)}~erit \bars{fere \textit{gestr.}} ratio %
\textit{(2)}~propemodum dici potest fore \textit{d} minorem quam %
\textit{(3)}~si \textit{d} sit mino %
\textit{(4)}~si numerus %
\textit{(a)}~\textit{rd} %
\textit{(b)}~\textit{d} exprimens rationem densitatis aeris ordinarii ad aerem inclusum, sit minor quam unitate minutus%
\textit{(5)}~propemodum dici potest %
\textit{(6)}~circiter dici %
\lbrack...\rbrack\ id est si %
\bars{\textit{d} \textit{erg.}} numerus %
\textit{(a)}~densit %
\textit{(b)}~exprimens rationem %
\lbrack...\rbrack\ pelliculae aqueae.\rbrack%
~\textit{L}%
}}
Sit verbi gratia $d = 10,$
seu aer ambiens decuplo densior incluso,%
\indexS{aer inclusus bullae}%
\indexS{aer ambiens}
fiet 9
\edtext{%
ad \textit{r} minus
quam 799 ad 800,%
}{%
\lemma{ad \textit{r}}%
\Bfootnote{%
\textit{(1)}~ut 799 %
\textit{(2)}~minus quam 799 ad 800,%
~\textit{L}%
}}
et \textit{r}
\edtext{%
majus quam
$\displaystyle\frac{7200}{799}$%
\rule[0mm]{0mm}{5.9mm}%
}{%
\lemma{majus quam}%
\Bfootnote{%
\textit{(1)}~\rule[0mm]{0mm}{5.0mm}$\displaystyle\frac{800 \cdot 9}{799}$ seu %
\textit{(2)}~$\displaystyle\frac{7200}{799}$%
~\textit{L}%
}}%
\edlabel{LH_37_03_123-124_124r_drittestufe-2}
\lbrack%
\textit{der Text bricht ab.}
\href{\Scan{LH_37_03_124v}}%
{124\,v\textsuperscript{o}}%
\rbrack%
\pend%
\vspace{1.50em}%
\normalsize
\pstart%
\noindent%
\lbrack%
\textit{Vierter Ansatz:}%
\rbrack\
\edtext{%
\edtext{%
\edlabel{LH_37_03_123-124_124v_viertestufe-1}%
\hebr{'}
Experimentis compertum est%
\indexS{experimentum}
gravitatem specificam aquae%
\indexS{gravitas specifica aquae}
circiter 800 vicibus continere
gravitatem specificam aeris%
\indexS{gravitas specifica aeris}%
}{%
\lemma{\hebr{'}}%
\Bfootnote{%
\textit{(1)}~Gravitas specifica %
\textit{(2)}~Experimentis compertum est gravitatem specificam\rule[0mm]{0mm}{3.6mm} %
\textit{(a)}~aeris in %
\textit{(b)}~aquae circiter %
\lbrack...\rbrack\ gravitatem specificam %
\textit{(aa)}~aquae %
\textit{(bb)}~aeris%
~\textit{L}%
}}
ordinarii seu ambientis.%
\indexS{aer ordinarius}%
\indexS{aer ambiens}%
}{%
\lemma{Experimentis \lbrack...\rbrack\ ambientis}%
\Cfootnote{%
Nicht ermittelt.%
}}
\edtext{%
Ponamus %
hunc \textit{d} vicibus
continere gravitatem specificam
seu densitatem aeris inclusi,%
\indexS{densitas aeris in bulla}
et pondus%
}{%
\lemma{Ponamus}%
\Bfootnote{%
\textit{(1)}~po %
\textit{(2)}~et %
\textit{(3)}~hunc \textit{d} vicibus %
\lbrack...\rbrack\ et pondus%
~\textit{L}%
}}
aeris%
\indexS{pondus aeris ordinarii}
\edtext{%
ordinarii
qui sit%
}{%
\lemma{ordinarii qui sit}%
\Bfootnote{%
\textit{erg.\;L}%
}}
paris cum bulla spatii%
\indexS{spatium bullae}%
\indexS{bulla vaporis}
\edtext{%
esse \textit{p}.
Bullae%
}{%
\lemma{esse \textit{p}.}%
\Bfootnote{%
\textit{(1)}~Ejus %
\textit{(2)}~Bullae%
~\textit{L}%
}}
centrum sit \textit{A},%
\indexS{centrum bullae}
sphaerae aeris inclusi radius \textit{AB},%
\indexS{sphaera aeris inclusi}
sphaerae totius bullae radius \textit{AC},%
\indexS{sphaera bullae}%
\indexS{radius bullae}
erit
\edlabel{LH_37_03_123-124_124v_liue-1}%
\edtext{}{%
\xxref{LH_37_03_123-124_124v_liue-1}{LH_37_03_123-124_124v_liue-2}%
\lemma{spatium}%
\Bfootnote{%
\textit{(1)}~bullae ut %
\textit{(a)}~rad %
\textit{(b)}~cub %
\textit{(2)}~aeris inclusi %
\lbrack...\rbrack\ ut cub.\,\textit{AC},%
~\textit{L}%
}}%
spatium aeris inclusi%
\indexS{spatium aeris inclusi}
ut cubus ab\;\textit{AB},
spatium totius bullae%
\indexS{spatium bullae}
ut cub.\,\textit{AC},%
\edlabel{LH_37_03_123-124_124v_liue-2}
spatium pelliculae aqueae%
\indexS{spatium pelliculae aqueae}%
\indexS{pellicula aquea bullarume}
ut differentia horum cuborum\lbrack;\rbrack%
\indexS{differentia cuborum}
quae
\edtext{%
sit
$\displaystyle\frac{1}{r}\,\text{cub.}\,AC,$%
}{%
\lemma{sit}%
\Bfootnote{%
\textit{(1)}~$\displaystyle\frac{r}{d}$ %
\textit{(2)}~$\displaystyle\frac{1}{r}\,\text{cub.}\,AC,$%
~\textit{L}%
}}
eritque
$\displaystyle\text{cub.}\,AB =
\edtext{%
\text{cub.}\,AC - \frac{1}{r}\,\text{cub.}\,AC.$%
\rule[0mm]{0mm}{5.8mm}
Pondus aeris inclusi%
\indexS{pondus aeris in bulla}%
\indexS{aer inclusus bullae}
erit
$\displaystyle\frac{1}{d}\frac{AB\,\text{cub.}}{AC\,\text{cub.}}p = \frac{1}{d} \cdot \frac{r - 1}{r}p.$%
\rule[0mm]{0mm}{5.6mm}
Pondus pelliculae%
\indexS{pondus pelliculae aqueae}%
\indexS{pellicula aquea bullarume}
erit
$\displaystyle\frac{800}{r}p.$
Ergo Pondus totius bullae%
\indexS{pondus bullae}
erit%
}{%
\lemma{$\displaystyle\text{cub.}\,AC - \frac{1}{r}\,\text{cub.}\,AC.$}%
\Bfootnote{%
\textit{(1)}~Pondus aqueae pelliculae %
\textit{(a)}~erit %
\textit{(b)}~ad pondus %
\textit{(aa)}~\textit{p} %
\textit{(bb)}~800\textit{pr} %
\textit{(c)}~in ratione co %
\textit{(d)}~\rule[0mm]{0mm}{4.0mm}ad %
\textit{(aa)}~pondus aeris %
\textit{(bb)}~pondus aeris inclusi erit %
\textit{(aaa)}~\textit{p} %
\textit{(bbb)}~\rule[0mm]{0mm}{5.0mm}$\displaystyle\frac{1}{d}\frac{AB\,\text{cub.}}{AC\,\text{cub.}}p$ seu $\displaystyle\frac{1}{d, \, 1 + \displaystyle\frac{1}{r}}p,$ %
\textit{(aaaa)}~ponderis %
\textit{(bbbb)}~pondus pelliculae erit %
\textit{(aaaaa)}~$\displaystyle\frac{800p}{r}$ %
\textit{(bbbbb)}~800 %
\textit{(2)}~Pondus aeris %
\lbrack...\rbrack\ erit
$\displaystyle\frac{800}{r}p.$ %
\bars{Ergo \textit{erg.}} Pondus totius bullae %
\bars{erit \textit{erg.}}%
\textit{L}%
}}
$\displaystyle(\frac{r - 1}{d} + 800)\frac{p}{r},$
quod debet esse minus
\ledpb%
quam \textit{p} %
pondus aeris paris spatii,%
\indexS{pondus aeris ambientis}
et
\edtext{%
fiet
$r - 1 + 800d$
minus quam \textit{rd},
adeoque%
}{%
\lemma{fiet}%
\Bfootnote{%
\textit{(1)}~$rp - 1p + 800dp$ %
\textit{(2)}~$r - 1 + 800d$ minus quam %
\textit{(a)}~\textit{rp} %
\textit{(b)}~\textit{rdp}, %
\textit{(c)}~\textit{rd}, %
\textit{(aa)}~seu %
\textit{(bb)}~adeoque%
~\textit{L}%
}}
\textit{r} majus quam
\edtext{%
$\displaystyle\frac{800d - 1}{d - 1}.$
Et pelliculae spatium%
\indexS{spatium pelliculae aqueae}%
\indexS{pellicula aquea bullarume}
designatum per~%
$\displaystyle\frac{1}{r}$
minus quam
$\displaystyle\frac{d - 1}{800d - 1}$
ipsius bullae.%
\indexS{bulla vaporis}%
\indexS{spatium bullae}
\edtext{%
\lbrack%
Itaque%
}{%
\lemma{\lbrack Itaque}%
\Cfootnote{%
Die eckige Klammer stammt von Leibniz.%
}}
si aer externus%
\indexS{aer externus}
sit decuplo%
}{%
\lemma{$\displaystyle\frac{800d - 1}{d - 1}.$}%
\Bfootnote{%
\rule[0.0mm]{0.0mm}{5.4mm}%
\textit{(1)}~Itaque si aer externus sit decuplo %
\textit{(2)}~Et pelliculae spatium %
\bars{designatum per $\displaystyle\frac{1}{r}$ \textit{erg.}} minus quam %
\lbrack...\rbrack\ sit decuplo%
~\textit{L}%
}}
densior incluso,%
\indexS{aer inclusus bullae}
fiet \textit{r} majus
quam%
\rule[0mm]{0mm}{5.4mm}
$7999 : 9,$
id
\edtext{%
est majus quam 889,%
}{%
\lemma{est majus quam}%
\Bfootnote{%
\textit{(1)}~890 %
\textit{(2)}~888 %
\textit{(3)}~889,%
~\textit{L}%
}}
eritque pelliculae spatium%
\indexS{spatium pelliculae aqueae}%
\indexS{pellicula aquea bullarume}
minus
\edtext{%
quam
$\displaystyle[\frac{1}{889}]$
spatii ipsius%
\indexS{spatium bullae}
\edtext{%
bullae.%
\rbrack%
}{%
\lemma{bullae.\rbrack\,}%
\Cfootnote{%
\rule[0.0mm]{0.0mm}{3.2mm}%
Die eckige Klammer stammt von Leibniz.%
}}
Unde ex regula cubica,%
\indexS{regula cubica}%
\rule[0.0mm]{0.0mm}{6.5mm}
ipsius quoque bullae crassities%
\indexS{crassities bullae}%
\indexS{bulla vaporis}
determinari potest.
Nam
$\displaystyle AB = AC \cdot \sqrt[3]{\vphantom{\displaystyle\frac{1}{r}}}\!(1 - \frac{1}{r})$
et
$\displaystyle AC - AB = AC, \: 1 - \sqrt[3]{\vphantom{\displaystyle\frac{1}{r}}}\!(1 - \frac{1}{r}).$
Jam
$\displaystyle 1 - \frac{1}{r}$
majus
\lbrack quam\rbrack\
$\displaystyle 1 - \frac{d - 1}{800d - 1},$
fit
$\displaystyle 1 - \frac{1}{r}$
majus quam
$\displaystyle\frac{800d - 1 -d + 1}{800d - 1},$
adeoque%
}{%
\lemma{quam}%
\Bfootnote{%
\bars{$\displaystyle\frac{1}{888}$ \textit{ändert Hrsg.}} %
\bars{spatii \textit{erg.}} ipsius\rule[0mm]{0mm}{3.6mm} bullae %
\textit{(1)}~; et ejus crassities %
\textit{(2)}~. Unde ex %
\lbrack...\rbrack\ determinari potest. %
\textit{(a)}~Fit enim cub.\,\textit{A} %
\textit{(b)}~Nam $\displaystyle AB = AC \cdot \sqrt[3]{\vphantom{\displaystyle\frac{1}{r}}}\!(1 - \frac{1}{r})$ et %
\bars{et \textit{streicht Hrsg.}} $\displaystyle AC - AB = AC, \: 1 - \sqrt[3]{\vphantom{\displaystyle\frac{1}{r}}}\!(1 - \frac{1}{r})$ %
\textit{(aa)}~, adeoque %
\textit{(bb)}~. Jam $\displaystyle 1 - \frac{1}{r}$ %
\textit{(aaa)}~= %
\textit{(bbb)}~majus \bars{quam \textit{erg. Hrsg.}} $\displaystyle 1 - \frac{d - 1}{800d - 1},$ fit $\displaystyle 1 - \frac{1}{r}$ %
\textit{(aaaa)}~= %
\textit{(bbbb)}~majus quam $\displaystyle\frac{800d - 1 -d + 1}{800d - 1},$ %
\textit{(aaaaa)}~adeoque %
\textit{(bbbbb)}~adeoque\rule[0mm]{0mm}{3.6mm}%
~\textit{L}%
}}
radius aeris inclusi%
\indexS{aer inclusus bullae}
major est quam%
\rule[0mm]{0mm}{5.6mm}
$\sqrt[3]{\vphantom{()}}(800d - d,\, :\, ,\, 800d - 1),$
ergo crassities pelliculae%
\indexS{crassities pelliculae aqueae}
\edlabel{LH_37_03_123-124_124v_verblasst-1}%
\edtext{}{%
\xxref{LH_37_03_123-124_124v_verblasst-1}{LH_37_03_123-124_124v_verblasst-2}%
\lemma{erit}%
\Bfootnote{%
\textit{(1)}~major %
\textit{(2)}~minor quam \textit{AC} %
\textit{(3)}~ad\rule[0mm]{0mm}{3.6mm} %
\textit{(a)}~crassi %
\textit{(b)}~\textit{AC} radium %
\lbrack...\rbrack\ ad unitatem. %
\textit{(aa)}~Unde %
\textit{(bb)}~Unde si %
\bars{si \textit{streicht Hrsg.}} aer bullae sit decuplo %
\bars{minor \textit{ändert Hrsg.}} aere ambiente erit %
\textit{(aaa)}~rat %
\textit{(bbb)}~crassitiei pelliculae %
\lbrack...\rbrack\ id est \dots\ \lbrack/\rbrack\ %
\textit{(aaaa)}~Ex hoc %
\textit{(bbbb)}~Si calculo simul intelligitur %
\textit{(cccc)}~Simili calculo aestimari potest%
~\textit{L}%
}}%
erit ad \textit{AC}
radium totius bullae%
\indexS{radius bullae}
in minore ratione quam est%
\rule[0mm]{0mm}{5.6mm}
$\displaystyle 1 - \sqrt[3]{\vphantom{\displaystyle\frac{1}{r}}}\frac{800d - d}{800d -1}$
ad unitatem.%
\indexS{unitas}
Unde
\edlabel{LH_37_03_123-124_124v_decuplo-1}%
si aer bullae%
\indexS{aer inclusus bullae}
sit decuplo
\lbrack minus densus\rbrack\
aere ambiente%
\indexS{aer ambiens}%
\edlabel{LH_37_03_123-124_124v_decuplo-2}
\ledpb%
erit crassitiei pelliculae ratio%
\indexS{crassities pelliculae aqueae}%
\indexS{pellicula aquea bullarume}
ad radium bullae%
\indexS{radius bullae}
in minore ratione quam est
\edtext{%
$1 - \sqrt[3]{\vphantom{()}}(890 : 899)$%
}{%
\lemma{\textit{Am Rand:}}%
\Afootnote{%
$\displaystyle
\efrac{8000-10}{8000-1}
=
\frac{890}{899}$}%
\lemma{$1 - \sqrt[3]{\vphantom{()}}(890 : 899)$}%
\Cfootnote{%
Leibniz setzt in seine Rechnung $d = 10$ ein (\refpassage{LH_37_03_123-124_124v_decuplo-1}{LH_37_03_123-124_124v_decuplo-2}).
Daraus ergibt sich eigentlich:
$\displaystyle 1 - \sqrt[3]{\frac{7.990}{7.999}}$.
Das falsche Ergebis ist bereits in der Nebenrechnung am Rand anzutreffen.%
}%
}
ad unitatem,%
\indexS{unitas}
id est
\dots
\pend%
\pstart%
Simili calculo%
\indexS{calculus}
aestimari potest%
\edlabel{LH_37_03_123-124_124v_verblasst-2}
quanta debeat
\edtext{%
esse magnitudo et crassities
sphaerae metallicae aere exhaustae,%
\indexS{magnitudo sphaerae aere exhaustae}%
\indexS{crassities sphaerae aere exhaustae}%
\indexS{sphaera metallica exhausta aere}
quae in aere nostro natare possit.%
\indexS{sphaera metallica natans in aere}%
\indexS{aer noster}%
}{%
\lemma{esse}%
\Bfootnote{%
\textit{(1)}~crassities %
\textit{(2)}~crassities sphaerae metallicae aere exhaus %
\textit{(3)}~magnitudo et crassities sphaerae metallicae %
\textit{(a)}~exhaustae, quae in %
\textit{(b)}~aere exhaustae, quae in aere %
\bars{nostro \textit{erg.}} natare possit.%
~\textit{L}%
}}
Pono specificam gravitatem metalli%
\indexS{gravitas specifica metalli}
esse ad gravitatem specificam aeris%
\indexS{gravitas specifica aeris}
ut \textit{m}
\edtext{%
ad 1,
et posito
recipientem esse sic satis exhaustum,%
\indexS{recipiens exhaustus}
fiet%
}{%
\lemma{ad\;1,}%
\Bfootnote{%
\textit{(1)}~fiet %
\textit{(2)}~et posito %
\lbrack...\rbrack\ exhaustum, fiet%
~\textit{L}%
}}
\textit{d} quasi numerus%
\indexS{numerus infinitus}
\edtext{%
infinitus
aere ambiente%
\indexS{aer ambiens}
incomparabiliter densiore existente
quam est aer in globo residens,%
\indexS{aer in globo exhausto}
et%
\rule[0.0mm]{0.0mm}{5.2mm}
$\displaystyle\frac{md - d}{md - 1}$
fiet
$\displaystyle\frac{m - 1}{m},$%
}{%
\lemma{infinitus}%
\Bfootnote{%
\bars{aere %
\textit{(1)}~\textlangle\textendash\textrangle\ %
\textit{(2)}~ambiente incomparabiliter %
\lbrack...\rbrack\ globo residens, \textit{erg.}} %
et %
\textit{(1)}~$\displaystyle\frac{800d - d}{800d - 1}$ fiet $\displaystyle=\: \frac{800 - 1}{800}$ %
\textit{(2)}~$\displaystyle\frac{md - d}{md - 1}$ fiet $\displaystyle\frac{m - 1}{m},$%
~\textit{L}%
}}
et posito sphaerae
\edtext{%
metallicae%
}{%
\lemma{metallicae}%
\Bfootnote{%
\textit{erg.\;L}%
}}
totius radium esse \textit{AC},
erit crassities%
\rule[0.0mm]{0.0mm}{4.0mm}
\edtext{%
metalli%
\indexS{crassities sphaerae metallicae}%
\indexS{crassities metalli}%
\indexS{metallum}
(quam vocamus \textit{K})%
}{%
\lemma{metalli}%
\Bfootnote{%
\textit{(1)}~\textit{AC} %
\textit{(2)}~(quam vocamus \textit{K})%
~\textit{L}%
}}
minor quam
\edtext{%
$\displaystyle AC,\, 1 - \sqrt[3]{\vphantom{\displaystyle\frac{1}{r}}}\frac{m - 1}{m}.$
\edlabel{LH_37_03_123-124-124v_nebenrechn-1}%
Itaque si metallum esset
5832 vicibus ponderosius aere%
\lbrack,\rbrack\
crassities metalli%
\indexS{crassities sphaerae metallicae}%
\indexS{crassities metalli}%
\indexS{metallum}
minor esse deberet
quam
\edtext{%
$\displaystyle AC,\, 1 - \frac{1}{18}\sqrt[3]{\vphantom{()}}5831,$
\ledpb%
seu
$\displaystyle AC,\, \frac{18000 - 17998{\scriptstyle\frac{1}{2}}}{18000} = AC,\, \frac{1{\scriptstyle\frac{1}{2}}}{18000} = \frac{\angles{{\scriptstyle\frac{3}{4}}}}{9000},\, AC = \frac{1}{12000} [AC].$%
\edlabel{LH_37_03_123-124-124v_nebenrechn-2}%
}{%
\lemma{\textit{Am unteren Rand:}}%
\Afootnote{%
$\sqrt[3]{\vphantom{()}}5831$
cadit inter 17.999 et 17.997.%
\textsuperscript{\lbrack a\rbrack} %
Nam \fbox{3}\:17.999 dat 5831.028
et \fbox{3}\:17.997 dat 5829.08[4].%
\textsuperscript{\lbrack b\rbrack}
Ergo sumamus
$\sqrt[3]{\vphantom{()}}5831 = 17.998[9].$%
\textsuperscript{\lbrack c\rbrack}%
\\[1.5ex]%
{\footnotesize%
\textsuperscript{\lbrack a\rbrack}~%
17.997.~%
\textit{(1)}~Nam illud proce %
\textit{(2)}~Nam \fbox{3}\:17.999 dat 5831.028%
~\textit{L}
\quad
\textsuperscript{\lbrack b\rbrack}~%
5829.086.%
~\textit{L~ändert Hrsg.}
\quad
\textsuperscript{\lbrack c\rbrack}~%
17.9985.%
~\textit{L~ändert Hrsg.}
}%
}}%
}{%
\lemma{$\displaystyle AC,\, 1 - \sqrt[3]{\vphantom{\displaystyle\frac{1}{r}}}\frac{m - 1}{m}.$}%
\Bfootnote{%
\textit{(1)}~Ergo \textit{K} minor quam $\displaystyle AC - AC \, \sqrt[3]{\vphantom{\displaystyle\frac{1}{r}}}\frac{m - 1}{m}$ seu $\displaystyle\sqrt[3]{\vphantom{\displaystyle\frac{1}{r}}}\frac{m - 1}{m}$ minor quam %
\textit{(a)}~$AC - K$ et %
\textit{(b)}~$\displaystyle 1 - \frac{K}{AC}$ et $\displaystyle\frac{m - 1}{m} = \text{cub.}\,(1 - \frac{K}{AC}),$ seu $\displaystyle m = \frac{1}{1 - \text{cub.}\,(1 - \displaystyle\frac{K}{AC})}$ %
\textit{(2)}~Itaque si metallum esset %
\textit{(a)}~8000 vicibus ponderosius aere, fieret crassities metalli minor quam $\displaystyle AC,\, 1 - \frac{1}{20} \sqrt[3]{\vphantom{j}} 7999$ %
\textit{(b)}~5832 vicibus ponderosius aere %
\textit{(aa)}~foret %
\textit{(bb)}~crassities metalli %
\lbrack...\rbrack\ $\displaystyle%
\frac{{\scriptstyle\langle\frac{3}{4}\rangle}}{9000},\, AC = \frac{1}{12000}$ %
\bars{\textit{AC} \textit{erg. Hrsg.}}.%
~\textit{L}%
}}
\pend%
\pstart%
Itaque
si ponamus
crassitiem metalli%
\indexS{crassities metalli}%
\indexS{metallum}
minimum
\edtext{%
esse duorum pollicum,%
\indexS{pollex}%
\rule[0.0mm]{0.0mm}{4.8mm}%
}{%
\lemma{esse}%
\Bfootnote{%
\textit{(1)}~pollic\textlangle em\textrangle, %
\textit{(2)}~duorum pollicum %
\textit{(3)}~duorum pollicum,%
~\textit{L}%
}}
erit radius%
\indexS{radius sphaerae metallicae}
\edtext{%
sphaerae 1500 pedum
et diameter 3000 pedum%
}{%
\lemma{sphaerae}%
\Bfootnote{%
\textit{(1)}~1000 pedum et diameter 2000 pedum %
\textit{(2)}~1500 pedum %
\lbrack...\rbrack\ 3000 pedum%
~\textit{L}%
}}
vel 600
\edtext{%
passuum geometricorum.%
\indexS{passus geometricus}%
}{%
\lemma{passuum geometricorum}%
\Cfootnote{%
Der \textit{passus geometricus} entsprach dem römischen Schritt, d.h. fünf Fuß.%
}}%
\edtext{}{%
\lemma{\textit{Zwischen den Zeilen:}}%
\Afootnote{%
Calculetur et quod
aucto globo,%
\indexS{globus aere exhaustus}
\textlangle ut\textrangle\ crescit pondus aeris incumbentis,%
\indexS{pondus aeris incumbentis}\indexS{aer incumbens}
ita et crassitie\lbrack s\rbrack\ auct\lbrack a\rbrack,%
\indexS{crassities globi aere exhausti}%
\textsuperscript{\lbrack a\rbrack}
sed augetur
sol\textlangle um\textrangle\
ut bases,
ergo ut superficies.%
\indexS{superficies globi aere exhausti}%
\\[1.5ex]%
{\footnotesize%
\textsuperscript{\lbrack a\rbrack}~%
crassitiem auctam \textit{L~ändert Hrsg.}
}%
}}
Itaque
quoties crassities metalli hujusmodi%
\indexS{crassities metalli}%
\indexS{metallum}
erit pollicum,%
\indexS{pollex}
toties diameter sphaerae%
\indexS{diameter sphaerae aere exhaustae}
minimum erit
\edtext{%
\lbrack300\rbrack%
}{%
\lemma{400}%
\Bfootnote{%
\textit{L\;ändert Hrsg.}%
}}
passuum%
\indexS{passus geometricus}
\edtext{%
geometricorum.%
}{%
\lemma{geometricorum.}%
\Bfootnote{%
\bars{minimum \textit{gestr.}}. Itaque%
~\textit{L}%
}}
Itaque
valde lapsi sunt
\edtext{%
Franciscus Lana,%
\indexN{\textls{Lana de Terzi}, Francesco 1631\textendash1687}%
}{%
\lemma{Franciscus Lana}%
\Cfootnote{%
Siehe F.\,\textsc{Lana\:de\:Terzi},
\textit{Prodromo}, erster Teil, cap.\;6 (Brescia 1670, S.\;52--61, bes. S.\;56\,f.);\cite{00069}
\textsc{Ders.}, 
\textit{Magisterium}, lib.\;VI, cap.\;3, artificium 46 (Bd.\;II, Brixen 1686, S.\;291\,F--294\,A, bes. S.\;292\,F--293\,B).\cite{01408}
Leibniz hatte den \textit{Prodromo} 1671 gelesen;
vgl. \textit{LSB} VIII,\;1 N.\;16.\cite{01103}%
}}
\edtext{%
Joh. Lomeierus%
\indexN{\textls{Lohmeier} (Lohmeierus, Lomeierus), Philipp 1648\textendash1680}%
}{%
\lemma{Joh. Lomeierus}%
\Cfootnote{%
Siehe 
P.\:\textsc{Lohmeier} und F.\,D.\:\textsc{Prescheur},
\textit{Exercitatio physica de artificio navigandi per aerem}
\lbrack Dissertation 1676\rbrack,
3. Ausg., Rinteln 1708.\cite{01422} %
Leibniz verwechselt im Text Philipp Lohmeier wohl mit seinem Vater Johann Georg.\indexN{\textls{Lohmeier}, Johann Georg, gest. 1680}%
}}
et alii viri docti,%
\indexS{vir doctus}
qui sphaeram%
\indexS{sphaera metallica exhausta aere}
\edtext{%
hujusmodi arte humana%
\indexS{ars humana}
facile fabricari posse
crediderunt,%
}{%
\lemma{hujusmodi}%
\Bfootnote{%
\textit{(1)}~fabricari %
\textit{(2)}~arte humana facile %
\textit{(a)}~posse crediderunt, %
\textit{(b)}~fabricari posse crediderunt,%
~\textit{L}%
}}
ut taceam
materiam tam puram%
\indexS{materia pura}
et uniformem%
\indexS{materia uniformis}
parari non posse,
ut tantae moli aeris incumbentis%
\indexS{moles aeris incumbentis}%
\indexS{aer incumbens}
ubique aequaliter resistat.%
\indexS{materia resistens ubique aequaliter}%
\edlabel{LH_37_03_123-124_124v_viertestufe-2}%
\edtext{}{%
\lemma{\textit{Am Rand:}}%
\Afootnote{%
Opus non minus quam imperator Sinensis.%
\indexS{imperator Sinensis}%
}}%
\pend%
\newpage%
\normalsize%
\pstart%
\noindent%
\lbrack\textit{In L,\;Bl.\:\href{\Scan{LH_37_03_123r}}{123\,r\textsuperscript{o}}\! überlieferte Nebenrechnungen%
:}\rbrack%
\pagelabel{LH_37_03_123-124_123r_nbnrchngn}%
\pend%
\vspace{1.5em}%
\pstart%
\centering%
\phantom{3339}199%
\hspace{40mm}%
\phantom{7239}195%
\\
\phantom{33}$\underline{\phantom{39}199\vphantom{j}}$%
\hspace{40mm}%
\phantom{72}$\underline{\phantom{39}195\vphantom{j}}$%
\\
\phantom{333}1791%
\hspace{40mm}%
\phantom{7239}975%
\\
\phantom{33}1791\phantom{1}%
\hspace{40mm}%
\phantom{72}1755\phantom{5}%
\\
\phantom{33}\underline{199\phantom{91}\vphantom{j}}%
\hspace{40mm}%
\phantom{72}\underline{195\phantom{95}\vphantom{j}}%
\\
\phantom{33}39601%
\hspace{40mm}%
\edtext{%
\phantom{72}37125%
}{\lemma{37125}%
\Cfootnote{Das Quadrat von 195 ist 38.025.}}%
\\
$\underline{\phantom{3339}199\vphantom{j}}$%
\hspace{40mm}%
$\underline{\phantom{7239}195\vphantom{j}}$%
\\
\phantom{3}356409%
\hspace{40mm}%
\phantom{3}185625%
\\
356409\phantom{9}%
\hspace{40mm}%
334125\phantom{5}%
\\
$\underline{39601\phantom{99}\vphantom{j}}$%
\hspace{40mm}%
$\underline{37125\phantom{95}\vphantom{j}}$%
\\
7880599%
\hspace{40mm}%
\edtext{%
7239375%
}{\lemma{7239375}%
\Cfootnote{Der Kubus von 195 ist 7.414.875.}}%
\pend%
\vspace{2.0em}%
\pstart%
\noindent%
\centering%
\phantom{31}195%
\hspace{35mm}%
\phantom{3}1190%
\hspace{35mm}%
\phantom{33}38315%
\\
$\phantom{3}\underline{\phantom{119}6\vphantom{j}}$%
\hspace{35mm}%
$\underline{37125\vphantom{j}}$%
\hspace{35mm}%
$\underline{\phantom{3338}198\vphantom{j}}$%
\\
\phantom{3}1170%
\hspace{35mm}%
38315%
\hspace{35mm}%
\phantom{3}306520%
\\
\phantom{31170}%
\hspace{35mm}%
\phantom{38315}
\hspace{35mm}%
344835\phantom{0}%
\\
\phantom{31170}%
\hspace{35mm}%
\phantom{38315}
\hspace{35mm}%
$\underline{38315\phantom{50}\vphantom{j}}$%
\\
\phantom{31170}%
\hspace{35mm}%
\phantom{38315}
\hspace{35mm}%
7586370%
\pend%
\vspace{1.0em}%
\pstart%
\centering%
\phantom{66}19%
\\
$\phantom{6}\underline{\phantom{6}19\vphantom{j}}$%
\\
\phantom{6}171%
\\
$\underline{\phantom{6}19\phantom{9}\vphantom{j}}$%
\\
\edtext{%
\phantom{6}351%
}{\lemma{351}%
\Cfootnote{Das Quadrat von 19 ist 361.}}%
\\
\phantom{66}19%
\\
3159
\\
$\underline{351\phantom{9}\vphantom{j}}$%
\\
\edtext{%
6669%
}{\lemma{6669}%
\Cfootnote{Der Kubus von 19 ist 6859.}}%
\pend%
\vspace{1.5em}%
\pstart%
\centering%
$\phantom{9 + 6}$
$\displaystyle 20 - 19.9 = \frac{1}{100}[\textit{!}]$
\hspace{20mm}%
\hspace{20mm}%
$9 + 6$
$\phantom{\displaystyle 20 - 19.9 = \frac{1}{100}[\textit{!}]}$
\pend%
\newpage%
\pstart%
\noindent%
\lbrack\textit{In L,\;Bl.\:\href{\Scan{LH_37_03_123v}}{123\,v\textsuperscript{o}}\! überlieferte Nebenrechnungen zum Textabschnitt \refpassage{LH_37_03_123-124_123v_nbnrchng-1}{LH_37_03_123-124_123v_nbnrchng-2}:}\rbrack%
\pagelabel{LH_37_03_123-124_123v_nbnrchngn}%
\pend%
\vspace{1.0em}%
\pstart%
\noindent%
\centering%
$\phantom{\cancel{0} \cancel{0} \cancel{0}} 1 0 \phantom{\cancel{0} 0 0 \ f \ 0.000}$%
\\%
$\phantom{\cancel{0}} \cancel{9} \cancel{9} \cancel{9} \cancel{1} 1 \phantom{0 0 \ f \ 0.000}$%
\\%
$\cancel{7} \cancel{2} \cancel{0} \cancel{0} \cancel{0} \cancel{0} 0 0 \ f \ 9.011$%
\\%
$\phantom{\cancel{0}} \cancel{7} \cancel{9} \cancel{9} \cancel{9} \cancel{9} 9 \phantom{0 \ f \ 0.000}$%
\\%
$\phantom{\cancel{0} \cancel{0}} \cancel{7} \cancel{9} \cancel{9} 9 \phantom{0 0 \ f \ 0.000}$%
\\%
$\phantom{\cancel{0} \cancel{0} \cancel{0}} \cancel{7} 7 \phantom{\cancel{0} 0 0 \ f \ 0.000}$
\pend%
\vspace{1.00em}%
\pstart%
\noindent%
\begin{edarrayl}
3 &
\edtext{%
[9]%
}{%
\lemma{3}%
\Bfootnote{%
\textit{L\;ändert Hrsg.}}}
        & 27 \\
4 &  16 & 64 \\
5 &  25 & 125 \\
6 &  36 & 216 \\
7 &     &
\end{edarrayl}%
\begin{edarrayl}
10 & 100 & 1000 & 1000
\end{edarrayl}%
\pend%
\vspace{0.5em}%
\pstart%
\noindent%
\centering%
$\displaystyle\frac{8000}{8000 - 10d}$%
\quad%
\quad%
\quad%
$\displaystyle\frac{20}{\sqrt[3]{\vphantom{()}}(8000 - 10d)} = \frac{AC}{AB}$
\pend%
\vspace{3.0em}%
\pstart%
\noindent%
\lbrack\textit{In L (Bl.\:\href{\Scan{LH_37_03_124r}}{124\,r\textsuperscript{o}}\!) überlieferte Nebenrechnungen:%
}\rbrack%
\pend%
\vspace{1.0em}%
\pstart%
\noindent%
\centering%
$\phantom{100300}1001$%
\\
$\phantom{100}\underline{\phantom{300}1001\vphantom{j}}$%
\\
$\phantom{100300}1001$%
\\
$\phantom{100}\underline{100100\phantom{1}\vphantom{j}}$%
\\
$\phantom{100}1002001$%
\\
$\phantom{100}\underline{\phantom{300}1001\vphantom{j}}$%
\\
$\phantom{100}1002001$%
\\
$\underline{100200100\phantom{1}\vphantom{j}}$%
\\
$1003003001$%
\\
$\underline{1000000000\vphantom{j}}$%
\\
$\phantom{100}3003001$%
\pend%
\newpage%
\pstart%
\noindent%
\lbrack\textit{In L,\;Bl.\:\href{\Scan{LH_37_03_124v}}{124\,v\textsuperscript{o}}\! überlieferte Nebenrechnungen zum Textabschnitt \refpassage{LH_37_03_123-124-124v_nebenrechn-1}{LH_37_03_123-124-124v_nebenrechn-2}:}\rbrack%
\pend%
\vspace{1.5em}%
\pstart%
\noindent%
\centering%
$a = 17990$%
\hspace{10mm}%
$b = 9$
\\
$\phantom{a = 17990}$%
\hspace{10mm}%
\underline{\phantom{b\,\ 8}9\vphantom{j}}%
\\
$\phantom{a = 17990}$%
\hspace{10mm}%
$\phantom{bl\ }81$
\\
$\phantom{a = 17990}$%
\hspace{10mm}%
$\underline{\phantom{b =\ }9\vphantom{j}}$%
\\
$\phantom{a = 17990}$%
\hspace{10mm}%
$\phantom{b\;}729$%
\pend%
\vspace{1.0em}%
\pstart%
\noindent%
\centering%
\edlabel{LH_37_03_123-124_124v_djg-1}%
$\text{cub.}\,(a + b) = a^3 + 3ab (a + b) + b^3$%
\edlabel{LH_37_03_123-124_124v_djg-2}%
\pend%
\vspace{1.0em}%
\pstart%
\noindent%
\centering%
$b^3 = 729$
\pend%
\vspace{1.0em}%
\pstart%
\noindent%
\centering%
$\phantom{0}4857300 = 3ab\phantom{0}$
\\
$\underline{\phantom{000000}9\vphantom{j}}\phantom{0 = 3abb}$
\\
$43715700 = 3abb$
\pend%
\vspace{1.5em}%
\pstart%
\noindent%
\centering%
$\phantom{3}aab =
\edtext{%
2012760900\phantom{0}%
}{%
\lemma{2012760900}%
\Cfootnote{%
$a^2b$ ist vielmehr $323.640.100 \cdot 9 = 2.912.760.900.$
Der Fehler wirkt sich auf den weiteren Verlauf der Nebenrechnungen ein.%
}}$%
\\
$\phantom{3aab =\ }\underline{\phantom{0000000}3\vphantom{j}\phantom{00}}\phantom{0}$
\\
$3aab = 6038282700\phantom{0}$
\\
$3abb\, = \underline{\phantom{000}4371570\vphantom{j}}0$
\\
$\phantom{3aab =\ }6042654270\phantom{0}$
\\
$\phantom{3aab =\ }%
\edtext{%
\underline{\phantom{000000000}9\vphantom{j}}\phantom{0}
}{%
\lemma{9}%
\Cfootnote{%
Der Summand ist eigentlich $b^3,$ d.h. 729
(vgl. \refpassage{LH_37_03_123-124_124v_djg-1}{LH_37_03_123-124_124v_djg-2}).
Der Fehler wirkt sich auf den weiteren Verlauf dieser Rechnung ein.%
}}$%
\\
$\phantom{3aab =\ }%
6042654279\phantom{0}$
\pend%
\vspace{1.5em}%
\pstart%
\noindent%
\centering%
\hspace{20mm}%
$\phantom{000}1799\phantom{0000000000}$%
\hspace{20mm}%
$\phantom{000}3236401\phantom{0000000}$%
\\
\hspace{20mm}%
$\phantom{00}\underline{\phantom{0}1799\vphantom{j}}\phantom{0000000000}$%
\hspace{20mm}%
$\phantom{00}\underline{\phantom{0000}1799\vphantom{j}}\phantom{0000000}$%
\\
\hspace{20mm}%
$\phantom{00}16191\phantom{0000000000}$%
\hspace{20mm}%
$\phantom{00}29127609\phantom{0000000}$%
\\
\hspace{20mm}%
$\phantom{0}16191\phantom{0}\phantom{0000000000}$%
\hspace{20mm}%
\edtext{%
$\phantom{0}20127609\phantom{0}\phantom{0000000}$%
}{%
\lemma{20127609}%
\Cfootnote{%
Vielmehr 29127609.
Der Fehler wirkt sich auf den weiteren Verlauf dieser Rechnung ein.
Der Kubus von 1.799 ist 5.822.285.399.%
}}%
\\
\hspace{20mm}%
$12593\phantom{00}\phantom{0000000000}$%
\hspace{20mm}%
$22654807\phantom{00}\phantom{0000000}$%
\\
\hspace{20mm}%
$\underline{1799\phantom{000}\vphantom{j}}\phantom{0000000000}$%
\hspace{20mm}%
$\underline{3236401\phantom{000}\vphantom{j}}\phantom{0000000}$%
\\
\hspace{20mm}%
$3236401$\phantom{0000000000}%
\hspace{20mm}%
$5732285399\phantom{0000000}$%
\pend%
\newpage%
\pstart%
\noindent%
\centering%
\edlabel{LH_37_03_123-124_124v_lasku-1}%
$\phantom{[000]}603828270000$%
\edtext{}{%
\xxref{LH_37_03_123-124_124v_lasku-1}{LH_37_03_123-124_124v_lasku-2}%
\lemma{603828270000 \lbrack...\rbrack\ 603871985709}%
\Cfootnote{%
Rechnungen unstimmig.%
}}%
\hspace{35mm}%
$\phantom{[}5732285399000\phantom{00]}$%
\\
$\phantom{[000]}\phantom{0000}43715700$%
\hspace{35mm}%
$\phantom{[000}6042654279\phantom{00]}$%
\\
$\phantom{[000]}\underline{\phantom{00000000000}9\vphantom{j}}$%
\hspace{35mm}%
$\edtext{%
[\underline{5738328053279\vphantom{j}}00]$%
}{%
\lemma{573833804427900}%
\Bfootnote{%
\textit{L\;ändert Hrsg.}%
}}%
\\
$\phantom{[000]}603871985709$%
\hspace{35mm}%
$\phantom{[000}603871985709\phantom{]}$%
\edlabel{LH_37_03_123-124_124v_lasku-2}%
\pend%
\vspace{1.5em}%
\pstart%
\noindent%
\centering%
$\phantom{\cancel{0}}\cancel{7}\cancel{7}[7]\phantom{\ f\ 000}$%
\edtext{}{%
\lemma{2}%
\Bfootnote{%
\textit{L\;ändert Hrsg.}}}%
\\
$\cancel{7}\cancel{9}\cancel{9}\cancel{9}\ f\ 888$%
\\
$\phantom{\cancel{0}}\cancel{9}\cancel{9}\cancel{9}\phantom{\ f\ 000}$%
\pend%
\vspace{1.0em}%
\pstart%
\noindent%
\centering%
$\phantom{00}16$%
\hspace{70mm}%
$\phantom{000}6$%
\\
$\phantom{00}\underline{16\vphantom{j}}$%
\hspace{70mm}%
$\phantom{00}\underline{\phantom{0}6\vphantom{j}}$%
\\
$\phantom{00}96$%
\hspace{70mm}%
$\phantom{00}36$%
\\
$\phantom{0}\underline{16\vphantom{j}\phantom{0}}$%
\hspace{70mm}%
$\phantom{0}\underline{\phantom{00}6\vphantom{j}}$%
\\
$\phantom{0}256$%
\hspace{70mm}%
$\phantom{0}216$%
\\
$\underline{\phantom{00}16\vphantom{j}}$%
\hspace{70mm}%
$\underline{\phantom{00}27\vphantom{j}}$%
\\
$1536$%
\hspace{70mm}%
$1512$%
\\
$\underline{256\phantom{0}\vphantom{j}}$%
\hspace{70mm}%
$\underline{432\phantom{0}\vphantom{j}}$%
\\
$4096$%
\hspace{70mm}%
$5832$%
\pend%
\vspace{1.0em}%
\pstart%
\noindent%
\centering%
$\phantom{00}\phantom{00}179$%
\hspace{25mm}%
$\phantom{00}32041$%
\\
$\phantom{00}\phantom{0}\underline{\phantom{0}179\vphantom{j}}$%
\hspace{25mm}%
$\phantom{00}\phantom{0}\underline{\phantom{0}179\vphantom{j}}$%
\\
$\phantom{00}\phantom{0}1611$%
\hspace{25mm}%
$\phantom{0}288369$%
\\
$\phantom{00}1253\phantom{0}$%
\hspace{25mm}%
$224287\phantom{0}$%
\\
$\phantom{00}\underline{179\phantom{00}\vphantom{j}}$%
\hspace{25mm}%
$\underline{32041\phantom{00}\vphantom{j}}$%
\\
$\phantom{00}32041$%
\hspace{25mm}%
5735339
\pend%
\vspace{1.0em}%
\pstart%
\noindent%
\centering%
$6000\ f\ 5\phantom{00}$%
\hspace{75mm}%
$200\ f\ 400$%
\\
$12\phantom{00\ f\ 00}$%
\hspace{75mm}%
$\phantom{0}5\phantom{00\ f\ 000}$%
\pend%
\vspace{0.5em}%
\pstart%
\noindent%
\centering%
$\phantom{\cancel{0}}\cancel{6}\phantom{\cancel{0}0\ f\ 000}$%
\hspace{20mm}%
$\phantom{00\ f\ 0}\phantom{00000}$%
\\
$\cancel{9}\cancel{0}\cancel{0}0\ f\ 750$%
\hspace{20mm}%
$\phantom{00\ f\ 0}18000$%
\\
$\cancel{1}\cancel{2}\cancel{2}\phantom{0\ f\ 000}$%
\hspace{20mm}%
$\phantom{00\ f\ 0}17997$%
\\
$\phantom{\cancel{0}}\cancel{1}\phantom{\cancel{0}0\ f\ 000}$%
\hspace{20mm}%
$\phantom{00\ f\ 0}\phantom{0000}3$%
\pend%
\newpage
\pstart%
\noindent%
\centering%
$\phantom{0000}\phantom{0000}17999$%
\hspace{35mm}%
$\phantom{0000}323964001$%
\\
$\phantom{0000}\phantom{000}\underline{\phantom{0}17999\vphantom{j}}$%
\hspace{35mm}%
$\phantom{0000}\phantom{000}\underline{\phantom{0}17999\vphantom{j}}$%
\\
$\phantom{0000}\phantom{000}161991$%
\hspace{35mm}%
$\phantom{000}2915676009$%
\\
$\phantom{0000}\phantom{00}161991\phantom{0}$%
\hspace{35mm}%
$\phantom{00}2915676009\phantom{0}$%
\\
$\phantom{0000}\phantom{0}161991\phantom{00}$%
\hspace{35mm}%
$\phantom{0}2915676009\phantom{00}$%
\\
$\phantom{0000}125993\phantom{000}$%
\hspace{35mm}%
$2267748007\phantom{000}$%
\\
$\phantom{0000}\underline{17999\phantom{0000}\vphantom{j}}$%
\hspace{35mm}%
$\underline{323964001\phantom{0000}\vphantom{j}}$%
\\
$\phantom{0000}323964001$%
\hspace{35mm}%
$5831028053999$%
\pend%
\vspace{3.00em}%
\pstart%
\noindent%
\centering%
$\phantom{00000}\phantom{00000}179999$%
\hspace{35mm}%
$\phantom{00000}32399640001$%
\\
$\phantom{00000}\phantom{0000}\underline{\phantom{0}179999\vphantom{j}}$%
\hspace{35mm}%
$\phantom{00000}\phantom{0000}\underline{\phantom{0}179999\vphantom{j}}$%
\\
$\phantom{00000}\phantom{0000}1619991$%
\hspace{35mm}%
$\phantom{0000}291596760009$%
\\
$\phantom{00000}\phantom{000}1619991\phantom{0}$%
\hspace{35mm}%
$\phantom{000}291596760009\phantom{0}$%
\\
$\phantom{00000}\phantom{00}1619991\phantom{00}$%
\hspace{35mm}%
$\phantom{00}291596760009\phantom{00}$%
\\
$\phantom{00000}\phantom{0}1619991\phantom{000}$%
\hspace{35mm}%
$\phantom{0}291596760009\phantom{000}$%
\\
$\phantom{00000}1259993\phantom{0000}$%
\hspace{35mm}%
$226797480007\phantom{0000}$%
\\
$\phantom{00000}\underline{179999\phantom{00000}\vphantom{j}}$
\hspace{33,8mm}%
$\underline{32399640001\phantom{00000}\vphantom{j}}$%
\\
$\phantom{00000}32399640001$%
\hspace{35mm}%
$5831902800539999$%
\pend%
\vspace{3.00em}%
\pstart%
\noindent%
\centering%
$\phantom{000000}\phantom{0000000}17999999$%
\hspace{35mm}%
$\phantom{0000000}%
\edtext{%
32399964000001$%
}{%
\lemma{32399964000001}%
\Cfootnote{%
Von Leibniz falsch übertragen.
Der Kubus von 17.999.999 ist
5,831999028000054e21.}}%
\\
$\phantom{000000}\phantom{000000}\underline{\phantom{0}17999999\vphantom{j}}$%
\hspace{35mm}%
$\phantom{0000000}\underline{\phantom{000000}17999999\vphantom{j}}$%
\\
$\phantom{000000}\phantom{000000}161999991$%
\hspace{35mm}%
$\phantom{000000}291599676000009$%
\\
$\phantom{000000}\phantom{00000}161999991\phantom{0}$%
\hspace{35mm}%
$\phantom{00000}291599676000009\phantom{0}$%
\\
$\phantom{000000}\phantom{0000}161999991\phantom{00}$%
\hspace{35mm}%
$\phantom{0000}291599676000009\phantom{00}$%
\\
$\phantom{000000}\phantom{000}161999991\phantom{000}$%
\hspace{35mm}%
$\phantom{000}291599676000009\phantom{000}$%
\\
$\phantom{000000}\phantom{00}161999991\phantom{0000}$%
\hspace{35mm}%
$\phantom{00}291599676000009\phantom{0000}$%
\\
$\phantom{000000}\phantom{0}161999991\phantom{00000}$%
\hspace{35mm}%
$\phantom{0}291599676000009\phantom{00000}$%
\\
$\phantom{000000}125999993\phantom{000000}$%
\hspace{35mm}%
$226799748000007\phantom{000000}$%
\\
$\phantom{000000}\underline{17999999\phantom{0000000}\vphantom{j}}$%
\hspace{35mm}%
$\underline{32399964000001\phantom{0000000}\vphantom{j}}$%
\\
$\phantom{000000}323999964000001$%
\hspace{35mm}%
$583199\phantom{000000000000000}$%
\pend%
\newpage%
\pstart%
\noindent%
\centering%
$\phantom{0000}\phantom{0000}17995$%
\hspace{35mm}%
$\phantom{0000}323820025$%
\\
$\phantom{0000}\phantom{000}\underline{\phantom{0}17995\vphantom{j}}$%
\hspace{35mm}%
$\phantom{000}\underline{\phantom{0000}\phantom{0}17995\vphantom{j}}$%
\\
$\phantom{0000}\phantom{0000}89975$%
\hspace{35mm}%
$\phantom{000}1619100125$%
\\
$\phantom{0000}\phantom{00}161955\phantom{0}$%
\hspace{35mm}%
$\phantom{00}2914380225\phantom{0}$%
\\
$\phantom{0000}\phantom{0}161955\phantom{00}$%
\hspace{35mm}%
$\phantom{0}2914380225\phantom{00}$%
\\
$\phantom{0000}125965\phantom{000}$%
\hspace{35mm}%
$2266740175\phantom{000}$%
\\
$\phantom{0000}\underline{17995\phantom{0000}\vphantom{j}}$%
\hspace{35mm}%
$\underline{323820025\phantom{0000}\vphantom{j}}$%
\\
$\phantom{0000}323820025$%
\hspace{35mm}%
$5827141349875$%
\pend%
\vspace{2.5em}%
\pstart%
\noindent%
\centering%
$\phantom{00000}78995$%
\\
$\phantom{0000}\underline{\phantom{0}78995\vphantom{j}}$%
\\
$\phantom{0000}394975$%
\\
$\phantom{000}710955\phantom{0}$%
\\
$\phantom{00}710955\phantom{00}$%
\\
$\phantom{0}631960\phantom{000}$%
\\
$\underline{552965\phantom{0000}\vphantom{j}}$%
\\
$6240210025$%
\pend%
\vspace{3.5em}%
\pstart%
\noindent%
\centering%
$\begin{array}{l}
5831000000000000%
\\
\cdot\phantom{83}\cdot\phantom{00}\cdot\phantom{00}\cdot\phantom{00}\cdot\phantom{00}\cdot%
\\
\hline%
\\
1%
\\
\hline%
\\
1%
\end{array}$%
\pagelabel{LH_37_03_123-124_124v_nbnrchngn-ende}%
\pend%
\selectlanguage{latin}%
\frenchspacing%
\normalsize%
